\documentclass[11pt,a4paper]{article}
\usepackage[margin=22mm]{geometry}
\usepackage{fontspec}
\setmainfont{Latin Modern Roman}
\setsansfont{Latin Modern Sans}
\setmonofont[Scale=0.85]{DejaVu Sans Mono}
\usepackage{babel}
\babelprovide[main,import]{czech}
\usepackage{amsmath,amssymb,amsthm,mathtools}
\renewcommand{\proofname}{Důkaz}
\usepackage{booktabs,longtable,array}
\usepackage{enumitem}
\usepackage{xcolor}
\usepackage{hyperref}
\usepackage{xurl}
\usepackage{mathrsfs}
\usepackage{fancyhdr}
\hypersetup{colorlinks=true,linkcolor=blue!45!black,urlcolor=blue!45!black,pdftitle={TWIST-J: oboustranná normová mez původního rozptylového stavu},pdfauthor={Analytická pracovní poznámka}}
\setlength{\parindent}{0pt}
\setlength{\parskip}{5pt}
\setlength{\emergencystretch}{3em}
\setlist{nosep,leftmargin=1.6em}
\pagestyle{fancy}
\fancyhf{}
\fancyhead[L]{\small\sffamily TWIST-J / oboustranná normová mez}
\fancyhead[R]{\small\sffamily NON-CANONICAL / L1}
\fancyfoot[C]{\small\thepage}
\setlength{\headheight}{15pt}
\newcommand{\norm}[1]{\left\lVert#1\right\rVert}
\newcommand{\abs}[1]{\left\lvert#1\right\rvert}
\newcommand{\ind}{\mathbf1}
\newcommand{\dd}{\,\mathrm d}
\newcommand{\Rea}{\operatorname{Re}}
\newcommand{\Ran}{\operatorname{Ran}}
\newcommand{\Ccal}{\mathcal C}
\newcommand{\Tcal}{\mathcal T}
\newcommand{\eps}{\varepsilon}
\newtheorem{theorem}{Tvrzení}[section]
\newtheorem{lemma}[theorem]{Pomocné tvrzení}
\theoremstyle{remark}\newtheorem{remark}[theorem]{Poznámka}
\numberwithin{equation}{section}

\usepackage{calc}
\usepackage{needspace}
\providecommand{\tightlist}{\setlength{\itemsep}{0pt}\setlength{\parskip}{0pt}}
\setcounter{tocdepth}{1}

\begin{document}

\thispagestyle{empty}
{\Large\sffamily\bfseries Oboustranná normová mez původního rozptylového stavu TWIST-J\par}
\vspace{3mm}
{\large Úplný spektrální účet a kvantitativní část komplementu\par}
\vspace{4mm}
Samostatný analytický dodatek, 29. října 2026\par
\textbf{NON-CANONICAL / L1. Interně přezkoumaný návrh důkazu.}

\section*{Hlavní nový výsledek}

Pro skutečný celý rozptýlený stav \(\chi=S\Psi\), s původním
normalizovaným zdrojem a pevným počátkem času, je sestrojen
pomocný celý volný stav \(\widehat\chi\). Zahrnuje oba vázané
stacionární prahy z obou stran i oba ostré konce původního
energetického okna. Pro každý čas \(n\ge0\) platí
\[
\boxed{\begin{aligned}
\|(U^n\Omega_+-U_0^n)\chi\|
&\le\|(U^n\Omega_+-U_0^n)\widehat\chi\|
 +\varepsilon_{\rm app},\\
\varepsilon_{\rm app}&<2\cdot10^{-9}.
\end{aligned}}
\]
Mez je odvozena přímo ve \textbf{skutečné celé stavové normě}.
Započítává celou spektrální hustotu druhého kontaktu, celý kruh
\(K\), obě pulzní větve, všech 26 štítků a původní paměť/reference.
Pomocný stav se nepřenormovává a vybraný jev se nemění.

Rozhodující přesná identita je
\[
\boxed{\|\chi(K,E)-\chi_\infty(K,E)\|^2
=\frac{32\pi f(E)^2}{|F_1(E)|^2}
\Re\frac1{F_2(K,E)+4/F_1(E)}.}
\]
Kladnost úplné spektrální míry a místní harmonická míra dávají
oboustranný integrovaný odhad i bez bodové meze celého
\(F_{2,\mathrm{comp}}\).

\section*{Další dokončené výsledky}

Celý komplement u vázaných prahů má kvalitativně doložené
holomorfní pokračování. Jeho koncové úseky mají výslovnou
normu \(2^{78}\). Celý hlavní záporný atomový komplement
má normu \(<2^{280}\) na polydisku s poloměry \(2^{-286}\),
včetně obou pulzních větví a pohyblivých koncových oprav.
Tato poslední mez ještě nezahrnuje ostatní atomy a kontinuum.

\vfill
\small
Číslo \(2\cdot10^{-9}\) je zaplacená aproximační chyba,
nikoli celá časová chyba. Tento dodatek neurčuje nový menší
globální čas. Přesný zbývající časový vstup je uveden v CT10--CT13.
Práce je oddělena od rezervace \#1816. Veřejný Canon v106 se
nemění a žádný podmíněný matematický závěr není vydáván za
uskutečněný experiment.
\newpage\normalsize
\tableofcontents
\clearpage

\section{Původní úloha, normalizace a návaznost}

\subsection{Přesně zachovaný zdroj a jev}

Používáme původní unitární model, úplné kontaktní funkce
a rozptylové předpoklady z [S1--S3]. Neprovádí se nová
postselekce ani změna zdrojové normalizace. Položme
\[
h=10^{-3},\quad E_*=8\pi/5,\quad I=[E_*-h,E_*+h],
\quad \beta=4\pi/5,\quad e=E-\beta,\quad g=2\pi/5.
\]
Původní normalizovaný objektový vstup je
\[
\psi_{\rm obj}
=[2h\rho_{\rm in}(H_1)]^{-1/2}
\mathbf1_I(H_1)\mathbf1_{[-h,h]}(P)a_0,
\quad a_0=\delta_{R=0,r=0}\,iB,
\quad \|\psi_{\rm obj}\|=1.
\]
Token je přesně v \(y=0\). Celková hybnost \(K\) proto
zůstává rovnoměrná na celém kruhu; energetická míra vstupu
je přesně \(\mathbf1_I(E)\,\mathrm dE/(2h)\).
Výstupní projektor je stále
\[
\begin{aligned}
Q={}&\mathbf1_{I_K}(K)\mathbf1_{|P|\le1/100}
\mathbf1_I\bigl(\Omega_b(P)+\beta+\varepsilon(k)\bigr)\Pi_b(P)\\
&\otimes P_+(k)\mathbf1_{\varepsilon'(k)>0}
\otimes I_{\rm label}\otimes|d_1\rangle\langle d_1|
\otimes I_{\rm ref}.
\end{aligned}
\]
Operátor \(I_{\rm label}\) zachovává všech 26 štítků.
Původní paměť může být v libovolné korelaci s referencí.
Místní dekompozice integračního kruhu a pomocné stavy slouží důkazu
regularity; žádný člen zdroje ani jevu se tím nevynechává.

\subsection{Úplný kontakt a rozptylový operátor}

Původní pulzní fáze a váhy jsou
\[
\varepsilon(k)=\arccos\frac{1-4\cos k}{5},\qquad
h_\sigma(k)=\frac{1-\sigma\varepsilon'(k)}2.
\]
Se správným vnějším předpisem \(E\mapsto E+i0\) platí
\[
F_1(E)=\int_{-\pi}^{\pi}f_B(P,E-\beta)\frac{\mathrm dP}{2\pi},
\]
\[
F_2(K,E)=\sum_{\sigma=\pm1}\int_{-\pi}^{\pi}
h_\sigma(K-P)f_A\bigl(P,E-\beta-\sigma\varepsilon(K-P)\bigr)
\frac{\mathrm dP}{2\pi},
\]
\[
T(K,E)=\frac1{4+F_1F_2}
\begin{pmatrix}-2F_2&-4\\4&-2F_1\end{pmatrix},\qquad
V(K,E)=T(K,E)e_1=\frac{(-2F_2,4)^T}{4+F_1F_2}.
\]
V druhém kontaktu je ponecháno celé objektové kontinuum,
všechny kontaktně viditelné póly a obě pulzní větve. Symbol
\(V(k)\) v místní maticové identitě pro pulz níže označuje
původní dvourozměrnou pulzní matici. Je odlišný od sloupce
\(V(K,E)\); argumenty obou symbolů jsou vždy uvedeny.


Výstupní okno je přesně
\[
I_K=[k_*-10^{-3},k_*+10^{-3}],\qquad
k_*=-\arccos\frac{9-5\sqrt5}{16}.
\]
Zdroj se na toto okno předem neomezuje. Jeho původní celé
K-rozložení má stále míru \(\mathrm dK/(2\pi)\).

\subsection{Které předchozí analytické výsledky se přebírají}

Tento dodatek navazuje na tři samostatné pracovní důkazy
[S4--S6]. Jejich stav je rovněž NON-CANONICAL/L1.
Značky BC, BS, CS, M, L označují úplný fyzikální soupis
a rychlostní důkazy v [S5]. Značky QA, QB, QC, FQ, BF
označují výslovné místní důkazy v [S6]. PT označuje
fyzikální hranový soupis v [S4]. Níže jsou nové důkazy
SP, CT, CG, HC a AM uvedeny v úplnosti.

Pro nový normový důkaz postačují tyto již doložené vstupy:
\begin{enumerate}
\item V původním okně jsou právě dvě vázané stacionární
konfigurace, obě s \(\sigma=-1\). Ve středu mají
\(r_b>1/32\), \(h_->3/8\) a společnou křivost
\(1/4<\Phi_{PP}<2\).
\item Původní objektový pól má holomorfní graf na
\(|P-P_0|<2^{-24}\), s \(|r_b|<64\),
\(|\omega_b''|<2^{20}\). Stacionární graf splňuje
\(|P_b'|,|E_b'|<8\) na \(|K-K_0|<2^{-58}\).
\item Celý první kontakt splňuje
\(\Re F_1>1/16\), \(|F_1|<3\) na každém komplexním
disku \(|E-E_0|<2^{-30}\), \(E_0\in I\).
\item Původní kontaktní spektrální mapa má
\(\mathcal Z^*\mathcal Z=\operatorname{diag}(\rho_1,\rho_2)\),
\(2\pi\rho_j=\Re F_j\), a skutečný operátor je
\(S=I+2\pi\mathcal ZT\mathcal Z^*\), podle [S1, §15].
\end{enumerate}
Převzaté obecné spektrální a rozptylové předpoklady se tím
nezesilují. Nové konstrukce nepřidávají fyzikální předpoklad
o rychlosti konvergence vlnových operátorů.

\subsection{Přehled přesného dosahu}

\begin{center}
\small
\begin{tabular}{>{\raggedright\arraybackslash}p{41mm}>{\raggedright\arraybackslash}p{49mm}>{\raggedright\arraybackslash}p{61mm}}
\toprule
Část důkazu & Nově doložená mez & Rozsah\\
\midrule
SP: úplný stav & \(2^{19}\eta^{3/4}\), \(\eta\le2^{-64}\) &
Obě strany obou vázaných prahů; celý druhý kontakt\\[3pt]
CT: prahy a konce I & \(\varepsilon_{\rm app}<2\cdot10^{-9}\) &
Pomocný stav a přesný návrat k původnímu výstupu\\[3pt]
CG: reálná geometrie & \(|t_P|>1/1600\), \(|D_o|>1/2000\) &
Všechny reálné hranové průchody v pásu BF\\[3pt]
HC: koncové úseky & \(|F_{2,\rm ann}|<2^{78}\) &
Dva přesné intervaly při pohyblivých koncích\\[3pt]
HC: celý komplement & Holomorfie na společném okolí &
Číselný poloměr celého okolí není přidělen\\[3pt]
AM: celý hlavní atom & \(|\mathcal A_{\rm comp}|<2^{280}\) &
Poloměry \(2^{-286}\); ostatní spektrum se nepřičítá\\
\bottomrule
\end{tabular}
\end{center}

\clearpage
\section{Oboustranná mez ve skutečné celé stavové normě}
\hypertarget{vuxfdsledek-oboustrannuxe1-norma-bez-bodovuxe9-meze-komplementu}{%
\subsection{Výsledek: oboustranná norma bez bodové meze
komplementu}\label{vuxfdsledek-oboustrannuxe1-norma-bez-bodovuxe9-meze-komplementu}}

Fixujme libovolný skutečný vázaný stacionární bod \((P_0,K_0,E_0)\) z
(BS25), s \(E_0\in I\). Holomorfní stacionární graf \(P_b(K)\),
\(E_b(K)\) je ten z (BF19),(BF20). Následující meze platí pro všechna
reálná

\[
|K-K_0|<2^{-60},\qquad 0<\eta\le2^{-64}.
\tag{SP1}
\]

Přípustný K-interval je zde větší než finální velmi malý polydisk BF2:
použije se již dříve dokázaný graf BF20, nikoli extrapolace koeficientu
A z BF3. Přípustné je také malé vybočení \(E_b(K)\) za konec I; potřebné
okolí prvního kontaktu stále obsahuje celý použitý interval.

Pro celý přenos a jeho limitní sloupec pišme

\[
V(K,E)=\frac{(-2F_2(K,E),4)^T}{4+F_1(E)F_2(K,E)},
\qquad V_\infty(E)=\binom{-2/F_1(E)}{0}.
\tag{SP2}
\]

Pak se bez předpokladu holomorfie nebo bodové meze
\(F_{2,\mathrm{comp}}\) dokáže

\[
\boxed{
\int_{E_b(K)-\eta}^{E_b(K)+\eta}
\|V(K,E)-V_\infty(E)\|_2^2\,dE
\le2^{56}\eta^{3/2}.}
\tag{SP3}
\]

Jde o \textbf{oboustranný integrovaný} výsledek pro skutečné úplné
okrajové hodnoty. Není to bodová Hölderova mez, holomorfní pokračování
celého komplementu ani nový globální čas.

Silnější fyzikálně vážená identita níže přenáší výsledek na celý původní
rozptýlený stav. Po spojení příslušných reálných K-intervalů okolo obou
vázaných křivek lze zvolit pomocný výstup \(\widetilde\chi\) tak, že

\[
\boxed{\|\chi-\widetilde\chi\|\le2^{18}\eta^{3/4},
\qquad \chi=S\Psi,}
\tag{SP4}
\]

a příspěvek návratu k původnímu stavu do odchozí bez přijímače časové
chyby je nejvýše

\[
\boxed{2^{19}\eta^{3/4}.}
\tag{SP5}
\]

Pro \(\eta=2^{-64}\) je (SP5) právě \(2^{-29}\). Původní normovaný
vstup, pevný počátek času a úplný vybraný jev se nepřepisují. Pomocný
výstup slouží pouze v trojúhelníkovém normovém účtu.

\hypertarget{pux16fvodnuxed-uxfaplnuxe1-kladnuxe1-spektruxe1lnuxed-muxedra}{%
\subsection{Původní úplná kladná spektrální
míra}\label{pux16fvodnuxed-uxfaplnuxe1-kladnuxe1-spektruxe1lnuxed-muxedra}}

Pro reálné K má původní úplný kontakt reprezentaci

\[
F_2(K,E+i y)=\int_{\mathbb T}
\frac{e^{-i(E+i y)}+e^{-i\theta}}
{e^{-i(E+i y)}-e^{-i\theta}}\,d\mu_K(\theta),
\quad y>0,\quad \mu_K(\mathbb T)=1.
\tag{SP6}
\]

Jde přesně o míru normovaného úplného druhého kontaktního vektoru. Je to
součet obou pulzních větví a integrál přes všechna P, s vahami
\(h_\sigma(K-P)\ge0\), \(h_++h_-=1\), a původní mírou dP/(\(2\pi\)).
Každá objektová spektrální míra v tomto součtu zahrnuje všechny
viditelné vázané i kontinuální složky.

Reálná část jednotlivého původního Carathéodoryho jádra je přesně

\[
\operatorname{Re}
\frac{e^{-i(E+i y)}+e^{-i\theta}}
{e^{-i(E+i y)}-e^{-i\theta}}
=\frac{\sinh y}{\cosh y-\cos(E-\theta)}.
\tag{SP7}
\]

Pro \(0<y\le1\), \(|E-\theta|\le y\), stačí \(\sinh y\ge y\),
\(\cosh y-1\le y^2\), \(1-\cos(E-\theta)\le(E-\theta)^2/2\). Proto

\[
\frac{\sinh y}{\cosh y-\cos(E-\theta)}\ge\frac1{2y}.
\tag{SP8}
\]

Například \(\cosh y-1\le y^2\) plyne integrací \(\cosh u<2\) pro
\(0\le u\le1\).

Na skutečné objektové vázané větvi je kladná podmíra

\[
d\mu_{K,b}
=\bigl[\Phi(\cdot,K)\bigr]_*
\left(h_-(K-P)\,r_b(P)\frac{dP}{2\pi}\right),
\quad
\Phi(P,K)=\beta+\omega_b(P)-\varepsilon(K-P).
\tag{SP9}
\]

Vzorec používá spektrální atom o hmotnosti \(r_b\). Jeho Carathéodoryho
reziduum je \(2ir_b\), ale jeho hmotnost není \(2r_b\). Proto se v (SP9)
žádný další faktor 2 nebo \(2\pi\) nevkládá.

\hypertarget{ux161irux161uxed-radiuxe1lnuxed-dolnuxed-mez-z-objektovuxe9ho-grafu}{%
\subsection{Širší radiální dolní mez z objektového
grafu}\label{ux161irux161uxed-radiuxe1lnuxed-dolnuxed-mez-z-objektovuxe9ho-grafu}}

Použijme již dokázané původní objektové okolí z BF:

\[
R_o=2^{-24},\quad |r_b(P)|<64,\quad
|\omega_b''(P)|<2^{20}
\qquad (|P-P_0|<R_o).
\tag{SP10}
\]

Cauchyho odhad na polovičním P-disku dává \(|r_b'|\le2^{31}\). Pro (SP1)
je podle BF20

\[
|P_b(K)-P_0|\le8|K-K_0|<2^{-57},\qquad
|E_b(K)-E_0|<2^{-57}.
\tag{SP11}
\]

Zvolme nyní reálný P-interval kolem skutečného stacionárního bodu
\(P_b(K)\), o poloviční šířce

\[
w=2^{-11}\sqrt\eta\le2^{-43}.
\tag{SP12}
\]

Celý interval leží v \(|P-P_0|<R_o/2\). Z (BS31) a (SP10),(SP11) plyne
na něm

\[
r_b(P)>\frac1{32}-2^{31}(2^{-57}+2^{-43})
=\frac1{32}-2^{-26}-2^{-12}>\frac1{64}.
\]

Na reálném pulzním pásu \(|\varepsilon''|\le1\). Vzhledem k původnímu
\(h_-(K_0-P_0)>3/8\) se pulzní váha na tomto intervalu změní o méně než
\((2^{-60}+2^{-57}+2^{-43})/2<1/8\). Tedy

\[
\boxed{r_b(P)>2^{-6},\qquad h_-(K-P)>2^{-2}.}
\tag{SP13}
\]

Tyto jsou skutečné fyzikální atomy: objektový graf zůstává uvnitř jediné
fyzikální mezery a pochází ze stejného jednoduchého původního pólu BF17.
Samotný analytický kandidát na opačném listě se nepoužívá.

Jelikož \(\Phi_P(P_b(K),K)=0\) a
\(|\Phi_{PP}|\le|\omega_b''|+|\varepsilon''|<2^{21}\), Taylorova formule
dává pro \(|p|\le w\)

\[
|\Phi(P_b(K)+p,K)-E_b(K)|
\le2^{20}p^2\le\eta/4.
\tag{SP14}
\]

V tomto kroku se nepotřebuje tvrzení o křivosti na celém intervalu w;
stačí odhad absolutní hodnoty druhé derivace a přesná stacionarita ve
středu.

Z (SP7)--(SP14) a kladnosti všech ostatních částí úplné míry dostáváme

\[
\operatorname{Re}F_2(K,E_b(K)+i\eta)
\ge\int_{P_b-w}^{P_b+w}
\frac{1}{2\eta}\,2^{-8}\frac{dP}{2\pi}
=\frac{2^{-20}}{\pi}\eta^{-1/2}
>2^{-22}\eta^{-1/2}.
\tag{SP15}
\]

Toto je dolní mez \textbf{celého} \(F_2\) na jeho původní vnější straně.
Interval w může být větší než pohyblivý Morseův řez BF25 s konci
\(x=\pm2^{-64}\). To není záměna \(B_{\mathrm{loc}}\) za celý zbytek:
(SP9) používá kladnou atomovou podmíru úplného \(F_2\), tedy podle
potřeby části ležící na obou stranách těchto umělých integračních konců.
Původní dekompozice BF3 zůstává přesná a beze změny. Žádný konec
intervalu ani reziduum nebyly z celého operátoru vynechány.

\hypertarget{omezenuxe1-analytickuxe1-reciprokuxe1-funkce-na-muxedstnuxedm-polodisku}{%
\subsection{Omezená analytická reciproká funkce na místním
polodisku}\label{omezenuxe1-analytickuxe1-reciprokuxe1-funkce-na-muxedstnuxedm-polodisku}}

Pro přehlednost v tomto oddílu označme \(z=E-E_b(K)\) a zvolme

\[
R=2^{-31},\qquad D_R=\{z:\ |z|<R,\ \operatorname{Im}z>0\}.
\tag{SP16}
\]

Podle (SP11) a FQ9 leží celý tento polodisk ve fyzikálně identifikovaném
okolí úplného prvního kontaktu: \(|E-E_0|<2^{-31}+2^{-57}<2^{-30}\). Na
něm tedy

\[
\operatorname{Re}F_1>1/16,\qquad |F_1|<3.
\]

Definujme

\[
\Gamma(E)=\frac4{F_1(E)},\qquad
G(K,E)=\frac1{F_2(K,E)+\Gamma(E)}.
\tag{SP17}
\]

V celém horním polodisku je původní \(F_2\) holomorfní se striktně
kladnou reálnou částí, zatímco

\[
\operatorname{Re}\Gamma
=\frac{4\operatorname{Re}F_1}{|F_1|^2}>\frac1{36}.
\]

Proto G nemá pól, je holomorfní, \(|G|\le36\), a jeho reálná část u je
kladná omezená harmonická funkce. Přesně platí

\[
|G|^2\le36\operatorname{Re}G,
\quad
0\le\operatorname{Re}G(K,E_b+i\eta)
\le|G(K,E_b+i\eta)|\le2^{22}\sqrt\eta.
\tag{SP18}
\]

Poslední mez je důsledkem (SP15). Žádný horní odhad \(F_2\) nebo jeho
komplementu se zde nepředpokládá.

\hypertarget{pux159esnuxe1-muxedstnuxed-harmonickuxe1-muxedra}{%
\subsection{Přesná místní harmonická
míra}\label{pux159esnuxe1-muxedstnuxed-harmonickuxe1-muxedra}}

Použití plného kruhového Poissonova vzorce bez znalosti \(F_1\) mimo
jeho doložené okolí by zde nebylo oprávněné. Pracujeme výhradně v
polodisku (SP16).

Poissonovo jádro jeho průměru, pro vnitřní bod \(i\eta\) a reálné
\(-R<x<R\), je

\[
\mathcal P_{R,\eta}(x)
=\frac\eta\pi\left(
\frac1{x^2+\eta^2}-\frac{R^2}{R^4+\eta^2x^2}\right).
\tag{SP19}
\]

Vzorec lze přímo odvodit biholomorfním zobrazením \(z\mapsto-(z/R+R/z)\)
polodisku na horní polorovinu. Bod \(i\eta\) se zobrazí na
\(i(R/\eta-\eta/R)\), zatímco na průměru je derivace reálné okrajové
souřadnice \((R^2-x^2)/(Rx^2)\). Dosazení do obyčejného polorovinného
jádra dává (SP19), v bodě x=0 spojitým pokračováním. Ekvivalentní zjevně
kladný tvar je

\[
\mathcal P_{R,\eta}(x)
=\frac{\eta(R^2-\eta^2)(R^2-x^2)}
{\pi(x^2+\eta^2)(R^4+\eta^2x^2)}.
\]

Pro \(|x|\le\eta\) a \(\eta\le R/2\) první zlomek v závorkách (SP19) je
nejméně \(1/(2\eta^2)\), druhý nejvýše \(1/R^2\le1/(4\eta^2)\). Tedy

\[
\mathcal P_{R,\eta}(x)\ge\frac1{4\pi\eta}
\qquad(|x|\le\eta).
\tag{SP20}
\]

Reprezentace omezené kladné harmonické funkce u obsahuje také nezáporný
příspěvek půlkružnice. Jeho vypuštění \textbf{z dolní meze} proto dává

\[
u(i\eta)\ge\int_{-\eta}^{\eta}
\mathcal P_{R,\eta}(x)u(x+i0)\,dx.
\]

Pro použití postačují téměř všude existující okrajové hodnoty omezeného
holomorfního G; případně se nerovnost dostane limitou menších polodisků
a Fatouovým lemmatem. Nezápornost na zbývající hranici je zajištěna
(SP18), ne odhadem operátoru mimo FQ.

Protože \(\eta\le2^{-64}<R/2\), (SP18),(SP20) dávají

\[
\boxed{
\int_{-\eta}^{\eta}
\operatorname{Re}G(K,E_b+x+i0)\,dx
\le4\pi\eta\operatorname{Re}G(K,E_b+i\eta)
<2^{26}\eta^{3/2}.}
\tag{SP21}
\]

\hypertarget{oboustrannuxfd-integrovanuxfd-pux159enos}{%
\subsection{Oboustranný integrovaný
přenos}\label{oboustrannuxfd-integrovanuxfd-pux159enos}}

Z úplné původní matice T plyne přesná identita

\[
V-V_\infty
=\binom{8/F_1^2}{4/F_1}\,G.
\tag{SP22}
\]

Na doloženém okolí \(|F_1|\ge1/16\), takže norma vektoru před G je menší
než \(8\cdot256+4\cdot16=2112<2^{12}\). Z (SP18),(SP21) tedy

\[
\int_{-\eta}^{\eta}|G(K,E_b+x)|^2dx
\le36\,2^{26}\eta^{3/2}<2^{32}\eta^{3/2},
\] \[
\int_{-\eta}^{\eta}\|V(K,E_b+x)-V_\infty(E_b+x)\|^2dx
<2^{24}2^{32}\eta^{3/2},
\]

což dokazuje (SP3). Odhad se týká obou stran x=0 současně, včetně
záporné strany, na níž samotný argument
\(\operatorname{Re}F_2\ge A/\sqrt x\) neplatí.

\hypertarget{pux159esnuxe1-vuxe1ux17eenuxe1-identita-v-celuxe9-fyzikuxe1lnuxed-vuxfdstupnuxed-normux11b}{%
\subsection{Přesná vážená identita v celé fyzikální výstupní
normě}\label{pux159esnuxe1-vuxe1ux17eenuxe1-identita-v-celuxe9-fyzikuxe1lnuxed-vuxfdstupnuxed-normux11b}}

Označme \(\mathcal Z_K(E):\mathbb C^2\to\mathfrak h_K(E)\) původní úplný
kontaktní spektrální operátor. Jeho sloupce z\_1,z\_2 leží v
ortogonálních uspořádáních a

\[
\mathcal Z_K(E)^*\mathcal Z_K(E)
=\operatorname{diag}(\rho_1(E),\rho_2(K,E)),
\qquad \rho_j=\frac{\operatorname{Re}F_j}{2\pi}.
\tag{SP23}
\]

Zde \(\rho_2\) označuje \textbf{celou} hustotu druhého kontaktu; není to
hustota vybraného vázaného nebo pravoběžného výstupu. Auditovaný úplný
stacionární vzorec je

\[
S(K,E)=I+2\pi\mathcal Z_K(E)T(K,E)\mathcal Z_K(E)^*.
\tag{SP24}
\]

Na původním vstupu platí přesně

\[
\mathcal Z_K(E)^*\Psi(K,E)=f(E)e_1,
\qquad f(E)=\sqrt{\rho_{\rm in}(E)/(2h)},\quad E\in I,
\tag{SP25}
\]

a nula mimo I; K má původní míru \(\mathrm dK/(2\pi)\). Ze zdrojové meze
\(|f|\le3/20\) nic nepřepočítáváme na podmíněně vybraný kanál. Tuto mez
lze také přímo zkontrolovat z úplného lichého vzorce:
\(\rho(P,E)<8/(21\pi)\) a \(I_P=[-h,h]\) dávají
\(f^2<4/(21\pi^2)<4/189<9/400\).

Definujme na původním I a použitých reálných FQ okolích pomocný úplný
unitární operátor

\[
S_\infty(K,E)=I-\frac{4\pi}{F_1(E)}|z_1(E)\rangle\langle z_1(E)|.
\tag{SP26}
\]

Na jednotkovém vektoru \(z_1/\sqrt{\rho_1}\) má násobitel
\(1-4\pi\rho_1/F_1=1-2\operatorname{Re}F_1/F_1 =-\overline{F_1}/F_1\),
jehož modul je jedna. Na ortogonálním doplňku je identita. Je tedy
skutečně unitární na celé volné energetické vrstvě; není to
renormalizovaný vybraný výstup. Mimo uvedený energetický obor jej lze
definovat jako identitu. Na \(\chi_\infty=S_\infty\Psi\) to nemá vliv,
protože původní zdroj má podporu v I.

Položme \(\chi_\infty=S_\infty\Psi\). Z (SP24),(SP25) máme

\[
\chi-\chi_\infty
=2\pi f\,\mathcal Z_K(E)(V-V_\infty).
\tag{SP27}
\]

Spektrální váhy, které by jinak při prahu mohly růst, se zde účtují
přesně. Dosazení (SP22),(SP23) dává

\[
\begin{aligned}
\|\mathcal Z_K(V-V_\infty)\|^2
&=\frac{16}{2\pi|F_1|^2}
\left(\operatorname{Re}F_2+
\frac{4\operatorname{Re}F_1}{|F_1|^2}\right)|G|^2\\
&=\frac{16}{2\pi|F_1|^2}\operatorname{Re}G.
\end{aligned}
\tag{SP28}
\]

Poslední rovnost je přesná identita
\(\operatorname{Re}(1/H)=\operatorname{Re}H/|H|^2\), \(H=F_2+4/F_1\).
Žádný horní odhad \(\rho_2\) není potřeba.

Proto ve skutečné úplné výstupní vrstvě

\[
\boxed{
\|\chi(K,E)-\chi_\infty(K,E)\|^2
=\frac{32\pi f(E)^2}{|F_1(E)|^2}
\operatorname{Re}G(K,E)
\le2^{10}\operatorname{Re}G(K,E).}
\tag{SP29}
\]

Číselně stačí \(32\cdot4\cdot(9/400)\cdot256<2^{10}\). Toto je
důležitější než nevažované (SP3): jde již o normu skutečného původního
výstupního vektoru.

\hypertarget{spojenuxed-obou-kux159ivek-a-pomocnuxe1-nuxe1hrada-bez-zmux11bny-zdroje}{%
\subsection{Spojení obou křivek a pomocná náhrada bez změny
zdroje}\label{spojenuxed-obou-kux159ivek-a-pomocnuxe1-nuxe1hrada-bez-zmux11bny-zdroje}}

Na původním energetickém okně jsou obě křivky BS25 hladké, monotónní v
příslušné energetické parametrizaci a oddělené v K. Jejich lokální grafy
BF20 se na překryvech shodují jednoznačností stacionárního pokračování.
Vezměme jejich K-obory, případně rozšířené uvnitř reálných okolí (SP1),
a označme jejich sjednocení \(\mathcal J\). Každému zde přísluší jediná
z uvažovaných vázaných energií \(E_b\)(K). Jejich oddělení z BS27 stačí,
aby tak malá rozšíření nesplynula. Nezávisle na jejich konkrétní délce
platí \(\int_\mathcal J dK/(2\pi)\le1\).

Zvolme libovolnou měřitelnou \(0\le\Theta(K,E)\le1\), jejíž podpora leží
v

\[
K\in\mathcal J,\qquad |E-E_b(K)|\le\eta,
\]

a definujme pomocný vektor

\[
\widetilde\chi=(1-\Theta)\chi+\Theta\chi_\infty.
\tag{SP30}
\]

Samotné \(\Theta\) může být hladce definováno i kousek za konci I.
Skutečná změna stavu má automaticky podporu v I, protože oba původní
vektory \(\chi\) a \(\chi_\infty\) jsou mimo I nulové. Tím se netvrdí,
že ostré konce zdroje již byly vyhlazeny; jejich samostatný normový účet
se tímto krokem nemění.

Integrace (SP29) nejprve přes I průnik \([E_b-\eta,E_b+\eta]\), poté
přes původní \textbf{celou} míru \(\mathrm dK/(2\pi)\), s (SP21), dává

\[
\|\chi-\widetilde\chi\|^2
\le2^{10}2^{26}\eta^{3/2}=2^{36}\eta^{3/2}.
\tag{SP31}
\]

To dokazuje (SP4), bez dělení velikostí \(\mathcal J\), bez nové
pravděpodobnostní normalizace a bez záměny celého výstupu za vybranou
událost. Mimo \(\mathcal J\) je pomocný vektor přesně původní. Jestliže
\(E_b\)(K) nepatrně leží mimo I, omezí se jen integrační interval
zdroje; odhad přes celý symetrický interval zůstává platnou horní mezí.

Oba S a \(S_\infty\) jsou unitární na celé volné vrstvě, takže
\(\|\widetilde\chi\|\le\|\Psi\|=1\); žádné přenormování pomocného
vektoru se neprovádí. Pro ostrou hodnotu \(\Theta\in\{0,1\}\) je dokonce
jeho norma přesně jedna. Pro časové derivace lze použít hladké
\(\Theta\) rovné jedné na \(|E-E_b|\le\eta/2\), nule na
\(|E-E_b|\ge\eta\), s libovolným z již doložených koncově plochých
přechodů. Odhad (SP31) zůstává beze změny.

Všechny vzorce jsou totožné na všech 26 nosičových štítcích a působí
jako identita na staré paměti a referenci. Proto (SP31) platí se stejnou
konstantou pro libovolný původní normalizovaný provázaný
paměťový/referenční vstup.

\hypertarget{pux159esnuxe9-muxedsto-v-ux10dasovuxe9m-uxfaux10dtu}{%
\subsection{Přesné místo v časovém
účtu}\label{pux159esnuxe9-muxedsto-v-ux10dasovuxe9m-uxfaux10dtu}}

Na volném a.c. podprostoru jsou \(U_0\) a \(\Omega_+\) izometrické. Pro
každé n proto přesně

\[
\begin{aligned}
\|(U^n\Omega_+-U_0^n)\chi\|
&\le\|(U^n\Omega_+-U_0^n)\widetilde\chi\|
+2\|\chi-\widetilde\chi\|\\
&\le\|(U^n\Omega_+-U_0^n)\widetilde\chi\|
+2^{19}\eta^{3/4}.
\end{aligned}
\tag{SP32}
\]

Tento krok je platný i pro nenormovaný pomocný vektor a je nezávislý na
čase n.~Proto pro konkrétní \(\eta=2^{-64}\)

\[
\boxed{
\|(U^n\Omega_+-U_0^n)\chi\|
\le\|(U^n\Omega_+-U_0^n)\widetilde\chi\|+2^{-29}.}
\tag{SP33}
\]

Vnitřek prahového pásku má nyní mimořádně jednoduchý \textbf{volný}
kontaktní profil. Z (SP23), (SP25) a (SP26) plyne

\[
\boxed{
\mathcal Z_K(E)^*\chi_\infty(K,E)
=-f(E)\frac{\overline{F_1(E)}}{F_1(E)}\,e_1.}
\tag{SP34}
\]

Druhá složka je přesně nula a první používá pouze již kontrolovaný první
kontakt a původní zdroj. Reálná okrajová funkce \(\overline{F_1(E)}\) má
na symetrickém FQ okolí holomorfní pokračování
\(F_1^\sharp(z)=\overline{F_1(\bar z)}\). Nenulový jmenovatel \(F_1\) je
již doložen. Tento profil lze tedy použít v přímém odchozím Cookově
postupu pro pomocný volný výstup.

Mimo vnitřek odřezaného pásku zůstává skutečný úplný profil
\(\mathcal Z_K(E)^*\chi\), včetně obou původních uspořádání. Jeho
derivace na přechodových a zbývajících intervalech, ostatní pravé
stacionární prahy a ostré konce I mají nadále vlastní analytický účet.
Rovnice (SP32) sama takový odhad neposkytuje.

Je také nutné odlišit dvě pomocné konstrukce. Pouhá náhrada
\(fV\mapsto f[(1-\Theta)V+\Theta V_\infty]\) v inverzní Duhamelově
kontaktní amplitudě není bez dalšího identitou pro
\(K_cJ^*U^m\Omega_+\widetilde\chi\). Její malá nevažovaná L² chyba by
sama nedokazovala malou celou stavovou chybu ani absolutní Cookův
součet. Zde se použil úplný stacionární stav (SP30) a přesná vážená
identita (SP28); díky nim je právě (SP32) oprávněný normový návrat k
původnímu experimentu.

\hypertarget{ostux159ejux161uxed-varianta-na-velmi-maluxe9m-pux16fvodnuxedm-bf-okoluxed}{%
\subsection{Ostřejší varianta na velmi malém původním BF
okolí}\label{ostux159ejux161uxed-varianta-na-velmi-maluxe9m-pux16fvodnuxedm-bf-okoluxed}}

Pro \(|K-K_0|<2^{-88}\), \(0<\eta\le2^{-152}\), lze přímo použít
Morseovu souřadnici BF24 a \(A(K,x^2)>2^{-8}\) pro
\(0\le x\le\sqrt\eta\). Párování obou kořenů x a −x v atomové míře dává

\[
\operatorname{Re}F_2(K,E_b+i\eta)
\ge\frac1\pi\int_0^{\sqrt\eta}
A(K,x^2)\frac{dx}{2\eta}
\ge\frac{2^{-9}}\pi\eta^{-1/2}>2^{-11}\eta^{-1/2}.
\tag{SP35}
\]

Stejný místní harmonický důkaz pak dává

\[
\int_{-\eta}^{\eta}\operatorname{Re}G(K,E_b+x)dx
<2^{15}\eta^{3/2},
\quad
\int_{-\eta}^{\eta}\|V-V_\infty\|^2dx<2^{45}\eta^{3/2},
\] \[
\|\chi-\widetilde\chi\|\le2^{13}\eta^{3/4},
\quad\text{časový návrat}\le2^{14}\eta^{3/4}.
\tag{SP36}
\]

Širší varianta (SP1)--(SP33) je pro konkrétní odříznutí vhodnější,
protože dovoluje \(\eta=2^{-64}\). Menší varianta slouží také jako
nezávislá kontrola původního koeficientu a míry BF.

\hypertarget{co-kladnost-nedokazuje-a-co-jeux161tux11b-musuxed-dodat-geometrie}{%
\subsection{Co kladnost nedokazuje a co ještě musí dodat
geometrie}\label{co-kladnost-nedokazuje-a-co-jeux161tux11b-musuxed-dodat-geometrie}}

Z úplné kladné míry a jedné pevné lokální odmocninové hustoty samotných
nelze odvodit jednotný bodový oboustranný modul ani společný holomorfní
poloměr \(B_{\mathrm{comp}}\). To není jen chybějící krok této metody.
Existuje jednoduchý abstraktní důvod:

Uvažujme pevnou kladnou podmíru \(\nu\) s hustotou \(c\theta^{-1/2}\) na
malém \(0<\theta<a\), s celkovou hmotností menší než jedna. Její
transformace má týž nenulový pravostranný záhyb v nule. V bodě E=−d je
její imaginární část záporná a řádu \(-d^{-1/2}\). Kladný atom v
\(\theta=-2d\) přispívá \(i m\cot(d/2)\). Volbou malé kladné hmotnosti
\(m=O(\sqrt d)\) lze zrušit původní imaginární část právě v E=−d;
zbývající hmotnost se umístí například v \(\theta=\pi\).

Přesně, označí-li \(H_\nu(-d)\) imaginární část transformace \(\nu\) a
\(M=1-\nu(\mathbb T)>0\), stačí

\[
m=-\frac{H_\nu(-d)+M\tan(d/2)}
{\cot(d/2)-\tan(d/2)}.
\]

Pro dostatečně malé d je \(0<m<M\), a míra
\(\nu+m\delta_{-2d}+(M-m)\delta_\pi\) je normovaná a kladná. Její
Carathéodoryho transformace má v −d hodnotu nula a stále má původní
lokální záhyb v nule. Při \(F_1\)=1 je tedy v −d \(V=(0,1)^T\), i když
limitní sloupec je \((-2,0)^T\). Atom v komplementu může být libovolně
blízko nule.

Tento příklad není tvrzením o skutečném TWIST-J spektru. Dokazuje pouze,
proč nelze potřebnou bodovou mez celého komplementu vyvodit ze samotné
kladnosti, normalizace a lokálního BF koeficientu. Oproti tomu
integrované (SP3),(SP31) platí i v přítomnosti takových imaginárních
kompenzací.

Skutečný TWIST-J má navíc fyzikální soupis BS/CS a společné analytické
lokální příspěvky QA, QB a QC. Ty mohou dále dodat holomorfii a číselnou
mez celé pravidelné periody, nebo postačující derivace jejího
kontaktního profilu mimo již zaplacené prahové pásky. Tato poznámka
jejich nevyčíslený komplement nezaměňuje za lokální
\(B_{\mathrm{loc}}\). Poskytuje alternativní úplný normový účet, který
při řešení časového problému nemusí čekat na bodovou mez komplementu
\textbf{uvnitř} těchto pásků. Žádný nový globální čas, hotový experiment
nebo nová hodnota původní pravděpodobnosti se tím samotným neprohlašuje.

\section{Konkrétní pomocný stav a přesné napojení na časovou řadu}
\hypertarget{vuxfdslovnuxfd-pux159echod-u-obou-vuxe1zanuxfdch-prahux16f}{%
\subsection{Výslovný přechod u obou vázaných
prahů}\label{vuxfdslovnuxfd-pux159echod-u-obou-vuxe1zanuxfdch-prahux16f}}

Zvolme přesně

\[
\eta=2^{-64},\qquad h=10^{-3},\qquad
\delta=h\,2^{-68},\qquad I=[E_*-h,E_*+h].
\tag{CT1}
\]

Položme \(p(u)=6u^5-15u^4+10u^3\) na \([0,1]\) a definujme funkci
\(\theta\) na reálné ose:

\[
\theta(t)=
\begin{cases}
1,&|t|\le1/2,\\
1-p(2|t|-1),&1/2<|t|<1,\\
0,&|t|\ge1.
\end{cases}
\tag{CT2}
\]

Je to skutečná \(C^2\) funkce, \(0\le\theta\le1\). Z
\(p'=30u^2(1-u)^2\le15/8\) a \(|p''|\le10/\sqrt3<6\) plyne

\[
|\theta'|<4,\qquad |\theta''|<24.
\tag{CT3}
\]

Na sjednocení obou K-oborů \(\mathcal J\) ze SP zvolme
\(\Theta(K,E)=\theta((E-E_b(K))/\eta)\); mimo ně \(\Theta\)=0. Hladkost
v K pro tento časový Fourierův argument není potřebná: pracuje se ve
vláknovém prostoru \(L^2(\mathbb T_K;\mathbb C^2)\). V energii máme

\[
|\partial_E\Theta|<4/\eta,\qquad
|\partial_E^2\Theta|<24/\eta^2.
\tag{CT4}
\]

Lokální grafy se na překryvech shodují a oba jejich K-obory jsou
oddělené. Rozšíření oborů u konců I lze vést uvnitř reálných okolí SP1.
\(\Theta\) smí být definováno také mimo I; původní zdroj i všechny níže
uvedené stavy jsou tam nulové. Tím nevzniká nový zdrojový jev.

\hypertarget{oba-pux16fvodnuxed-ostruxe9-konce-zdroje}{%
\subsection{Oba původní ostré konce
zdroje}\label{oba-pux16fvodnuxed-ostruxe9-konce-zdroje}}

Nechť \(q_\delta(E)\) je nula mimo I, jedna mimo oba koncové proužky
šířky \(\delta\), a na každém proužku roste od nuly k jedné podle
\(r(u)=3u^2-2u^3\), orientovaného od příslušného konce I. Funkce
\(q_\delta\) je \(C^1\), má integrabilní druhou slabou derivaci a
\(0\le q_\delta\le1\). Původní zdroj má přesnou energetickou míru
\(\mathbf1_I(E)\,\mathrm dE/(2h)\), takže

\[
\|(1-q_\delta(H_0))\Psi\|^2
=\frac1{2h}\,2\delta\int_0^1(1-3u^2+2u^3)^2\,\mathrm du
=\frac{13\delta}{35h}.
\tag{CT5}
\]

Jde o normový účet původního normalizovaného zdroje, nikoli o jeho nové
přenormování. Oba ostré konce jsou započteny.

\hypertarget{jeden-celuxfd-pomocnuxfd-vuxfdstup-a-jeho-cena}{%
\subsection{Jeden celý pomocný výstup a jeho
cena}\label{jeden-celuxfd-pomocnuxfd-vuxfdstup-a-jeho-cena}}

Označme jako ve SP \(\chi=S\Psi\), \(\chi_\infty=S_\infty\Psi\) a zvolme

\[
\widehat\chi(K,E)=q_\delta(E)
\bigl[(1-\Theta(K,E))\chi(K,E)
 +\Theta(K,E)\chi_\infty(K,E)\bigr].
\tag{CT6}
\]

Tento vektor má normu nejvýše jedna. Protože S je unitární a na původním
bezpečném okně I komutuje s původní energií, z CT5 a SP31 plyne

\[
\|\chi-\widehat\chi\|
\le\sqrt{\frac{13\delta}{35h}}+2^{18}\eta^{3/4}.
\]

Pro každý nezáporný čas n jsou oba evoluční členy izometrické. Proto
návrat od pomocného stavu k původnímu celému výstupu stojí nejvýše

\[
\begin{aligned}
\varepsilon_{\rm app}
&:=2\sqrt{\frac{13\delta}{35h}}+2^{19}\eta^{3/4}\\
&=2\sqrt{\frac{13}{35}}\,2^{-34}+2^{-29}
<2^{-33}+2^{-29}=17\,2^{-33}<2\cdot10^{-9}.
\end{aligned}
\tag{CT7}
\]

Poslední porovnání je přesná racionální nerovnost
\(17\cdot10^9<2\cdot2^{33}\). Takto jsou současně a bez nové normalizace
zaplaceny oba vázané prahy z obou stran i oba původní ostré konce I.
Hodnota CT7 stále není celou časovou chybou.

\hypertarget{skuteux10dnuxfd-uxfaplnuxfd-volnuxfd-kontaktnuxed-profil}{%
\subsection{Skutečný úplný volný kontaktní
profil}\label{skuteux10dnuxfd-uxfaplnuxfd-volnuxfd-kontaktnuxed-profil}}

Položme \(b_\chi=\mathcal Z^*\chi\). Z původního úplného S a
\(\mathcal Z^*\mathcal Z =\operatorname{diag}(\operatorname{Re}F_1,\operatorname{Re}F_2)/(2\pi)\)
plyne přesně

\[
b_\chi(K,E)=f(E)
\left[e_1+
\begin{pmatrix}\operatorname{Re}F_1&0\\0&\operatorname{Re}F_2\end{pmatrix}
V(K,E)\right],\qquad E\in I.
\tag{CT8}
\]

Tedy plný profil pomocného stavu je

\[
\boxed{
\widehat b(K,E)=q_\delta(E)f(E)
\left\{(1-\Theta)
\left[e_1+\operatorname{diag}(\operatorname{Re}F_1,
\operatorname{Re}F_2)V\right]
-\Theta\frac{\overline{F_1}}{F_1}e_1\right\}.}
\tag{CT9}
\]

Na \(|E-E_b(K)|\le\eta/2\) zůstává pouze
\(-q_\delta f\,\overline{F_1}/F_1\) v prvním kontaktu. Na přechodech a
mimo pásky jsou zachovány celé \(F_2\) i celý přenosový sloupec. Druhá
složka v CT8 není nahrazena hustotou vybraného výstupu. Všechny identity
mají původní míru dE, nikoli \(\mathrm dE/(2\pi)\).

Profil CT9 je omezená měřitelná funkce. Z
\(\operatorname{Re}F_1\ge1/16\), \(|F_1|\le3\),
\(\operatorname{Re}F_2\ge0\) plyne \(|T_{11}|\le32\) a
\(|\operatorname{Re}F_2\,T_{21}|\le64\). První složka hranaté závorky
CT8 má tedy modul nejvýše 97, druhá 64. S \(|f|\le3/20\),
\(0\le q_\delta,\Theta\le1\) a \(|\overline{F_1}/F_1|=1\) dostáváme
\(\|\widehat b(K,E)\|_2<32\). Konečnost vážené normy CT10 z této samotné
bodové meze neplyne a bude dále samostatným číselným vstupem.

\hypertarget{pux159esnux11b-vymezenuxfd-zbuxfdvajuxedcuxed-ux10duxedselnuxfd-vstup-pro-novuxfd-ux10das}{%
\subsection{Přesně vymezený zbývající číselný vstup pro nový
čas}\label{pux159esnux11b-vymezenuxfd-zbuxfdvajuxedcuxed-ux10duxedselnuxfd-vstup-pro-novuxfd-ux10das}}

Nechť \(\mathcal H_K=L^2(\mathbb T_K,\mathrm dK/(2\pi);\mathbb C^2)\).
Pro skutečný profil CT9 definujme

\[
c_m(K)=\int_I e^{-imE}\widehat b(K,E)\,\mathrm dE,
\qquad
\mathcal B^2=\sum_{m\ge1}m^{3/2}\|c_m\|_{\mathcal H_K}^2.
\tag{CT10}
\]

Původní spektrální věta s mírou dE dává přesně
\(c_m=J^*U_0^m\widehat\chi\); nepřidává se faktor \(2\pi\) ani posun m.
Pro \(W_m=U^{-m}U_0^m\) platí konečná teleskopická identita

\[
(W_{m+1}-W_m)\widehat\chi
=-U^{-(m+1)}U_0J K_c\,c_m.
\]

Po sečtení od m=n, použití definice \(\Omega_+\) a vynásobení \(U^n\) má
tedy odpovídající Cookova řada přesný tvar

\[
(U^n\Omega_+-U_0^n)\widehat\chi
=-\sum_{m=n}^{\infty}U^{n-m-1}U_0J K_c\,c_m,
\]

jakmile je doložena níže uvedená absolutní konvergence.

Předpokládejme, že je dodáno \textbf{konkrétní doložené číslo} \(B\) s
\(\mathcal B\le B<\infty\). Potom pro \(n\ge2\) přímá odchozí Cookova
řada a \(\|U-U_0\|\le2\), v přesném kontaktním rozkladu
\(U-U_0=U_0JK_cJ^*\), dávají

\[
\begin{aligned}
\|(U^n\Omega_+-U_0^n)\widehat\chi\|
&\le2\sum_{m=n}^{\infty}\|c_m\|_{\mathcal H_K}\\
&\le2B\left(\sum_{m=n}^{\infty}m^{-3/2}\right)^{1/2}
\le\frac{2\sqrt2\,B}{(n-1)^{1/4}}.
\end{aligned}
\tag{CT11}
\]

Zde lze použít přesnější normu původního \(K_c\); bezpečná obecná mez 2
však stačí. Absolutní konvergence této řady je důsledkem dodaného CT10,
nikoli skrytým předpokladem při součtu.

Spojení s CT7 dává podmíněný nový odhad pro původní holý výstup

\[
\boxed{
\varepsilon_{+,\rm bare}(n)
<2\cdot10^{-9}+
\frac{2\sqrt2\,B}{(n-1)^{1/4}}.}
\tag{CT12}
\]

Pro libovolný cíl \(\varepsilon>2\cdot10^{-9}\) by stačil konkrétní čas

\[
n\ge\max\left\{2,\,
1+\left\lceil
\left(\frac{2\sqrt2\,B}{\varepsilon-2\cdot10^{-9}}\right)^4
\right\rceil\right\}.
\tag{CT13}
\]

Toto je přesná postačující podmínka, nikoli zde nově přidělené číselné
B. K jeho doložení musí být započteny přechody CT2, zbytek původního
profilu a další skutečné kontinuální prahy. Vyhlazení konců zruší
nesčitatelné holé krajní členy 1/m, ale jeho derivační cenu je třeba
zahrnout do B. Samotná kvalitativní holomorfie komplementu nebo již
známá lokální mez \(B_{\mathrm{loc}}\) tuto úlohu neplní.

Pokud se použije dřívější úplný konečný časový návrh, jeho časy a
pravděpodobnostní závěry se nemění. CT12--CT13 přesně určují jinou cestu
k novému časovému odhadu a jeden přesně vymezený postačující číselný
údaj v této cestě. Příchozí čas, kalibrace čtení a aparátové chyby mají
své původní samostatné účty. Žádný z nich se nezaměňuje s CT7.

\clearpage
\section{Fyzikální oddělení a všechny reálné hranové průchody}
Pišme

\[
g=2\pi/5,\quad e=E-4\pi/5,\quad
I_e=[4\pi/5-10^{-3},4\pi/5+10^{-3}],\quad
a=\arccos(-3/5),\quad b=\pi-a.
\tag{CG1}
\]

Pulz má kladnou fázi \(\varepsilon(k)\in[0,a]\) a velikost rychlosti

\[
W(\varepsilon)^2=
\frac{3+5\cos\varepsilon}{5(1+\cos\varepsilon)}.
\tag{CG2}
\]

Kruhová vzdálenost hybností se vždy počítá modulo \(2\pi\).

\hypertarget{ostux159ejux161uxed-fyzikuxe1lnuxed-obor-vuxe1zanuxe9ho-prahu}{%
\subsection{Ostřejší fyzikální obor vázaného
prahu}\label{ostux159ejux161uxed-fyzikuxe1lnuxed-obor-vuxe1zanuxe9ho-prahu}}

Použijme parametrizaci BS4--BS6:

\[
c=\cos\delta,\quad s=\sin\delta,\quad
v=c-4sz,\quad \mathscr G(z,\delta)=cU(z)-sV(z)=0,
\]

\[
U=(3+5C)z^2-5(1-C),\qquad
V=h(1+z^2)+4z-4z^3,
\quad C=\cos g,\quad h=5\sin g.
\tag{CG3}
\]

Již platí \(1/4<\delta_*<1/3\), \(1<z<5/4\), \(v'<-4\),
\(z'=V/(c\mathscr G_z)>0\), \(|v'|<64/9\), a

\[
309/1000<C<31/100,\qquad 19/4<h<119/25.
\tag{CG4}
\]

\hypertarget{dvux11b-pomocnuxe9-meze-pro-parametrizaci}{%
\subsubsection{Dvě pomocné meze pro
parametrizaci}\label{dvux11b-pomocnuxe9-meze-pro-parametrizaci}}

Na celém BS boxu máme

\[
z>9/8,\qquad z'>7/10.
\tag{CG5}
\]

První mez plyne z testu při \(z=9/8\): \(U<591/256\), \(V>2449/256\),
tedy \(U/V<591/2449<1/4<\tan\delta\). Funkce \(\mathscr G/z^2\) je
striktně rostoucí; její jediný kladný kořen proto leží napravo od testu.

Pro druhou mez je V konkávní na \([1,5/4]\), a proto

\[
V>\min\{19/2,599/64\}=599/64.
\]

Současně \(V'>-7/2\), \(3+5C<91/20\), takže

\[
0<c\mathscr G_z\le\mathscr G_z
<91/8+7/6=301/24.
\]

Tedy \(z'>1797/2408>7/10\). Z (CG5) vychází

\[
-v'=s+4cz+4sz'
>\frac{19}{5}s+4cz
>\frac{19}{5}\frac{95}{384}
  +4\frac{17}{18}\frac98
=\frac{1993}{384}>5.
\tag{CG6}
\]

Na každém stacionárním bodě proto \(w=\sqrt{1-v^2}/(-v')<1/5\).

\hypertarget{dolnuxed-mez-delta}{%
\subsubsection{Dolní mez delta}\label{dolnuxed-mez-delta}}

Z \(W=w<1/5\) plyne přesně \(\cos\varepsilon<-7/12\). Kdyby
\(\delta_*\le3/10\), bylo by \(\varepsilon\le3\pi/5+301/1000\). Pro
\(q=301/1000\) dávají střídavé Taylorovy meze \(\cos q<239/250\),
\(\sin q<3/10\). Proto

\[
\cos(3\pi/5+q)
=-C\cos q-\sin g\sin q
>-\frac{31}{100}\frac{239}{250}
 -\frac{119}{125}\frac3{10}
=-\frac{14549}{25000}>-\frac7{12},
\]

což je spor. Tedy

\[
\delta_*>3/10.
\tag{CG7}
\]

\hypertarget{hornuxed-mez-delta}{%
\subsubsection{Horní mez delta}\label{hornuxed-mez-delta}}

Testujme \(\delta_0=13/40\), \(z_0=121/100\). Taylorovy meze dávají

\[
\cos\delta_0>3031/3200,\qquad
\sin\delta_0<639/2000.
\]

Při tomto z je

\[
U>6398669/2000000,\qquad V<1185359/125000.
\]

Odtud přesně

\[
\mathscr G(z_0,\delta_0)
>\frac{18945367}{32000000000}>0.
\tag{CG8}
\]

Jedinečný fyzikální parametr proto splňuje \(z(\delta_0)<121/100\), a
tudíž

\[
v(\delta_0)>
\frac{3031}{3200}
-4\frac{639}{2000}\frac{121}{100}
=-\frac35+\frac{323}{400000}.
\tag{CG9}
\]

Proto \(\sin p>4/5\), a BS mez \(-v'<64/9\) dává \(w(\delta_0)>9/80\).

Je-li \(\delta_0\) za energeticky přípustným koncem
\(\delta_{\max}(e)\), je tvrzení \(\delta_*<\delta_0\) okamžité. Jinak
je při testu \(\varepsilon\ge3\pi/5+81/250\). Při \(q=81/250\) platí

\[
C\cos q+\sin g\sin q
>\frac{309}{1000}\left(1-\frac{q^2}{2}\right)
 +\frac{19}{20}\left(q-\frac{q^3}{6}\right)
=\frac{371997462}{625000000}
>\frac{3759}{6319}.
\]

Poslední nerovnost má kladný křížový rozdíl 1276962378. Z (CG2) je
\(\cos\varepsilon<-3759/6319\) právě podmínkou \(W<9/80\). Test tedy
leží na kladné straně rozdílu BS23. Jeho jediný průchod z minus do plus
nastává před testem:

\[
\boxed{3/10<\delta_*<13/40.}
\tag{CG10}
\]

\hypertarget{meze-v-rychlosti-a-celkovuxe9-hybnosti}{%
\subsubsection{Meze v, rychlosti a celkové
hybnosti}\label{meze-v-rychlosti-a-celkovuxe9-hybnosti}}

Na \(\delta>3/10\) test \(z=7/6\) dává

\[
U<395/144,\qquad V>4117/432,
\qquad U/V<1185/4117<3/10<\tan\delta,
\]

tedy \(z>7/6\). Jelikož \(c<24/25\) a \(s>591/2000\), máme

\[
v<\frac{24}{25}-4\frac{591}{2000}\frac76
=-\frac{419}{1000}<-\frac25.
\]

Z monotonicity v a (CG9) současně \(v>-3/5\). Navíc

\[
-v'>\frac{19}{5}\frac{591}{2000}
   +4\frac{17}{18}\frac76
=\frac{1493183}{270000}>\frac{11}{2}.
\]

Tedy

\[
\boxed{-3/5<v<-2/5,\qquad
 w<\frac{2\sqrt{21}}{55}<\frac16.}
\tag{CG11}
\]

Poslední ostrá nerovnost je přesně \(3024<3025\). Pro původní kladné P
je \(P=\arccos(-v)\), \(K_b=P-\kappa\). Střídavé meze poskytují

\[
\cos(37/40)>3/5,
\qquad \cos(7/40)<69/70,
\qquad \cos(6/5)<2/5.
\]

Z \(w<1/6\) a přesné pulzní formule je \(\cos\kappa>69/70\), takže
\(\kappa<7/40\). Proto

\[
\boxed{3/4<K_b(e)<6/5,\qquad K_\pm(e)=\pm K_b(e).}
\tag{CG12}
\]

Existuje i výslovná rezerva pro úzké BF okolí. Taylorův polynom šestého
řádu v 37/40 dává

\[
\cos(37/40)-3/5>
\frac{5371601591}{2949120000000}>10^{-3}.
\]

Z jednotkové Lipschitzovy meze kosinu plyne \(P>37/40+10^{-3}\). Dále
\(\cos(6/5)<229/625\), a tedy \(6/5-\arccos(2/5)>21/625>3/100\). Můžeme
tak zesílit (CG12) na

\[
\boxed{3/4+10^{-3}<K_b(e)<6/5-3/100.}
\tag{CG13}
\]

\hypertarget{oddux11blenuxed-od-vux161ech-ostatnuxedch-fyzikuxe1lnuxedch-konfiguracuxed}{%
\subsection{Oddělení od všech ostatních fyzikálních
konfigurací}\label{oddux11blenuxed-od-vux161ech-ostatnuxedch-fyzikuxe1lnuxedch-konfiguracuxed}}

Označme \(\varphi=(1+\sqrt5)/2\), \(v_a=(5\sqrt5-11)/4\),
\(p_a=\arccos v_a\), \(p_b=\pi-p_a\). Z \(559/250<\sqrt5<56/25\) plyne

\[
809/500<\varphi<81/50,\qquad
9/200<v_a<1/20,\qquad 3/2<p_a<31/20.
\tag{CG14}
\]

Poslední dvojici dokazují přímo \(\cos(3/2)>3231/46080>1/20\) a
\(\cos(31/20)<150721/3840000<9/200\).

\hypertarget{setkuxe1nuxed-puxf3lu-s-hranou}{%
\subsubsection{Setkání pólu s
hranou}\label{setkuxe1nuxed-puxf3lu-s-hranou}}

První čtveřice QA má \(K=\pm p_a\pm\kappa_a(e)\), kde

\[
\cos\kappa_a(e_0)=(5\varphi-3)/8,
\quad e_0=4\pi/5,
\quad
|\cos\kappa_a(e)-\cos\kappa_a(e_0)|\le1/800.
\]

Tedy \(127/200<\cos\kappa_a<511/800\), a Taylorovy testy
\(\cos(9/10)<127/200\), \(\cos(17/20)>511/800\) dávají

\[
17/20<\kappa_a<9/10.
\]

Její kladné absolutní polohy jsou proto

\[
3/5<p_a-\kappa_a<7/10,
\qquad 47/20<p_a+\kappa_a<49/20.
\tag{CG15}
\]

Druhá čtveřice má \(K=\pm p_b\pm\kappa_b\),
\(\cos\kappa_b(e_0)=(7-5\varphi)/8\). Z téže změny kosinu vychází

\[
-111/800<\cos\kappa_b<-27/200,
\qquad 3/2<\kappa_b<7/4.
\]

Horní úhlovou mez potvrzuje \(\cos(7/4)<-863/6144<-111/800\). Také
\(3/2<p_b<7/4\), tedy

\[
|p_b-\kappa_b|<1/4,
\qquad 3<p_b+\kappa_b<7/2.
\tag{CG16}
\]

Poslední součet má po redukci modulo \(2\pi\) absolutní hodnotu větší
než \(5/2\). S (CG12) dostáváme od \textbf{všech osmi} QA křivek
kruhovou vzdálenost větší než \(1/20\).

\hypertarget{kuux17eelovuxe9-konce-a-pulznuxed-banky}{%
\subsubsection{Kuželové konce a pulzní
banky}\label{kuux17eelovuxe9-konce-a-pulznuxed-banky}}

U přístupného kužele QC je kladné \(K_q=\pi-\kappa_q\),

\[
\cos K_q=-(1+5\cos e)/4,
\qquad \cos K_q(e_0)=(5\varphi-2)/8.
\]

Odtud \(\cos K_q>19/25\), zatímco \(\cos(18/25)<293749/390625<19/25\).
Proto

\[
0<K_q<18/25.
\tag{CG17}
\]

Obě kuželové křivky jsou od BF vzdáleny o více než \(3/100\).

Pro hranové setkání pulzních bank QB má kladná redukovaná poloha
\(K_h=\pi-\arccos(4+5\cos e)\), tedy

\[
\cos K_h=-4-5\cos e,
\qquad 1/25<\cos K_h<11/200.
\]

Proto \(3/2<K_h<\pi/2\), a obě bankové křivky jsou od BF vzdáleny o více
než \(3/10\).

\hypertarget{bux11bux17enuxe9-kontinuuxe1lnuxed-prahy-a-druhuxfd-vuxe1zanuxfd-pruxe1h}{%
\subsubsection{Běžné kontinuální prahy a druhý vázaný
práh}\label{bux11bux17enuxe9-kontinuuxe1lnuxed-prahy-a-druhuxfd-vuxe1zanuxfd-pruxe1h}}

Již dokázaná kontinuální poloha splňuje \(0<K_c(e)<12/25\). Z (CG12) se
její vzdálenost od BF zesiluje na \(27/100\). Oba samotné vázané prahy
mají kruhovou vzdálenost \(2K_b>3/2\), protože \(2K_b<12/5<\pi\).

Celý výsledek je následující. Meze v posledním sloupci platí pro
libovolný bod pásu \(3/4<|K|<6/5\), vyjma posledního řádku, který
porovnává dva skutečné BF body při stejné energii.

\Needspace{12\baselineskip}
\begin{longtable}[]{@{}
  >{\raggedright\arraybackslash}p{(\columnwidth - 4\tabcolsep) * \real{0.3000}}
  >{\raggedright\arraybackslash}p{(\columnwidth - 4\tabcolsep) * \real{0.3000}}
  >{\raggedleft\arraybackslash}p{(\columnwidth - 4\tabcolsep) * \real{0.4000}}@{}}
\toprule\noalign{}
\begin{minipage}[b]{\linewidth}\raggedright
Fyzikální konfigurace
\end{minipage} & \begin{minipage}[b]{\linewidth}\raggedright
Její redukované polohy
\end{minipage} & \begin{minipage}[b]{\linewidth}\raggedleft
Odstup od BF pásu/body
\end{minipage} \\
\midrule\noalign{}
\endhead
\bottomrule\noalign{}
\endlastfoot
Běžné kontinuální prahy & \(\pm K_c\), \(0<K_c<12/25\) & \(>27/100\) \\
První pólové a hranové čtveřice & (CG15) a jejich zápory & \(>1/20\) \\
Druhá pólové a hranové čtveřice & (CG16) modulo \(2\pi\) & \(>1/20\) \\
Přístupný kužel & \(\pm K_q\), \(0<K_q<18/25\) & \(>3/100\) \\
Hranové setkání bank & \(\pm K_h\), \(3/2<K_h<\pi/2\) & \(>3/10\) \\
Druhý vázaný práh & \(-K_b\) proti \(K_b\) & \(>3/2\) \\
\end{longtable}

\[
\operatorname{dist}_{\mathbb T}
\bigl(K_b(e),\{K_c,K_{QA},K_q,K_h,-K_b\}\bigr)>3/100.
\tag{CG18}
\]

U kontinuálních, QA, QC a QB křivek jsou v množině zahrnuty obě
znaménkové kopie; z dvojice BF je zde zahrnut pouze opačný bod. Nejsou
přidávány nepřístupné algebraické kořeny jiných vrstev.

\hypertarget{vux161echny-reuxe1lnuxe9-hranovuxe9-prux16fchody-majuxed-jednotnou-rychlostnuxed-mez}{%
\subsection{Všechny reálné hranové průchody mají jednotnou rychlostní
mez}\label{vux161echny-reuxe1lnuxe9-hranovuxe9-prux16fchody-majuxed-jednotnou-rychlostnuxed-mez}}

\hypertarget{jedinuxe1-potenciuxe1lnux11b-pomaluxe1-rodina}{%
\subsubsection{Jediná potenciálně pomalá
rodina}\label{jedinuxe1-potenciuxe1lnux11b-pomaluxe1-rodina}}

Pro nízkou \(t=-4\) hranu a kladný pulz parametrizujme
\(\omega\in[e-a,b]\). Absolutní objektová rychlost je

\[
u(\omega)=
\sqrt{\frac{\cos\omega-3/5}{1+\cos\omega}},
\qquad D(\omega)=W(e-\omega)-u(\omega).
\tag{CG19}
\]

Funkce u klesá a W(e−\(\omega\)) roste. Jejich jediný průsečík je
dřívější kontinuální práh \(\omega_c(e)\), s \(\omega_c<21/50\). V celém
přípustném intervalu je \(\omega>149/1000\), \(\sin\omega>1/8\),
\(u<1/2\). Proto

\[
-u'(\omega)
=\frac{(4/5)\sin\omega}{u(1+\cos\omega)^2}>1/20,
\quad D'>1/20.
\tag{CG20}
\]

Derivační formule se používají na otevřeném intervalu; v krajích, kde u
nebo W mizí, se níže uvedené závěry berou spojitým limitem.

Kdyby v některém bodě \(|D|\le10^{-3}\), pak
\(|\omega-\omega_c|\le1/50\) a \(\omega<11/25\). Po celém spojujícím
intervalu tedy \(\cos\omega>9/10\), \(u>1/3\), \(W>3/10\), \(uW>1/10\).
Pro protisměrné momenty \(K_-(\omega)=p(\omega)-\kappa(\omega)\) platí

\[
K_-'=\frac1u-\frac1W=\frac{D}{uW}.
\]

Monotonicita D proto dává

\[
|K_-(\omega)-K_c(e)|
<10\,(10^{-3})\,(20\cdot10^{-3})=1/5000.
\]

To odporuje (CG18) pro K v BF pásu. V protisměrné konfiguraci je tudíž
relativní rychlost větší než \(10^{-3}\). Pro souhlasné momenty je
rychlostní rozdíl součtem u+W. Je-li \(|D|>10^{-3}\), je tento součet
ještě větší; v opačném případě předchozí argument dává \(u>1/3\). Také v
této konfiguraci je součet větší než \(10^{-3}\).

\hypertarget{vux161echny-ostatnuxed-hranovuxe9-rodiny}{%
\subsubsection{Všechny ostatní hranové
rodiny}\label{vux161echny-ostatnuxed-hranovuxe9-rodiny}}

Přípustné zbývající hrany mají \(t=+4\) a pulzní fázi
\(0\le\varepsilon<\pi/2\). Pro kladný pulz to plyne z
\(\varepsilon\le e-a<\pi/2\). Pro záporný pulz použijme

\[
a>7\pi/10+10^{-3}.
\tag{CG21}
\]

Nerovnost (CG21) je elementární: \(\sin(\pi/5)<59/100\), takže
\(\cos(7\pi/10+10^{-3})=-\sin(\pi/5+10^{-3})>-591/1000>-3/5\). Proto
\(\varepsilon\le2\pi-a-e<\pi/2\). Pulz má tedy \(W>\sqrt{3/5}>3/4\),
kdežto každá jednostranná objektová hranová rychlost je nejvýše
\(1/\sqrt5<1/2\). Jejich relativní rychlost je větší než \(1/4\).

Celkově na každém skutečném reálném hranovém průchodu v BF pásu

\[
\boxed{|\omega_{\rm edge}'(P)-\sigma\varepsilon'(K-P)|>1/1000.}
\tag{CG22}
\]

\hypertarget{od-relativnuxed-rychlosti-k-pux16fvodnuxedmu-diskriminantu}{%
\subsubsection{Od relativní rychlosti k původnímu
diskriminantu}\label{od-relativnuxed-rychlosti-k-pux16fvodnuxedmu-diskriminantu}}

Pro nízkou hranu již máme \(\sin\omega>1/8\). Na vysoké hraně v BF pásu
platí

\[
|\sin\omega|>1/1000.
\tag{CG23}
\]

Pro kladný pulz je tato mez okamžitá, protože \(a\le\omega\le e<\pi\) a
\(\sin e>1/2\). U záporného pulzu by její selhání vyžadovalo
\(\omega=\pi+\nu\), \(|\nu|<1/500\). Přesná hranová identita dává

\[
|\sin((P-\pi)/2)|=\sqrt5|\sin(\nu/2)|,
\quad |P-\pi|<3|\nu|.
\]

Pulzní hybnost se při přechodu od \(\omega\) k \(\pi\) změní nejvýše o
\((4/3)|\nu|\), neboť celá příslušná pulzní fáze leží v \((0,\pi/2)\).
Celkové K by tedy bylo vzdáleno méně než \((13/3)|\nu|<13/1500<3/100\)
od kuželové křivky, což odporuje (CG18). Tím je (CG23) dokázáno.

Na hraně \(t=\cos P-5\cos\omega_\sigma\) má přesně derivaci

\[
\partial_Pt
=5\sin\omega\,[\sigma\varepsilon'(K-P)-\omega_{\rm edge}'(P)].
\]

Nízká hrana tedy dává \(|\partial_Pt|>1/1600\), vysoká dokonce
\(|\partial_Pt|>1/800\). Jednotná mez je

\[
\boxed{|\partial_Pt|>1/1600.}
\tag{CG24}
\]

\hypertarget{kontaktnuxed-jmenovatel-na-vux161ech-hranuxe1ch}{%
\subsection{Kontaktní jmenovatel na všech
hranách}\label{kontaktnuxed-jmenovatel-na-vux161ech-hranuxe1ch}}

Položme \(x=e^{-i(\omega-g)}\). Na nízké hraně \(q=-1\) původní
nezkrácený jmenovatel přesně dává

\[
|D_o|=8\cos((\omega-g)/2)\sin(g/2)
                 |\sin(\omega-g/2)|.
\tag{CG25}
\]

Prefaktor je větší než \(16/5\). Jestliže by \(|\omega-g/2|\le10^{-3}\),
potom podél hranové větve od \(\omega\) k g/2 platí
\(|v-v_a|\le5\cdot10^{-3}\), \(|P_\omega|<6\). Pulz v této malé oblasti
má \(W>1/2\), tedy \(|\kappa_\omega|<2\). Bod K by byl do
\(8\cdot10^{-3}\) od první QA čtveřice, v rozporu s odstupem \(1/20\).
Proto \(|\omega-g/2|>10^{-3}\). Tento úhlový rozdíl má modul menší než
jedna; z \(\sin y\ge5y/6\) na \([0,1]\) vychází \(|D_o|>1/375\).

Na vysoké hraně \(q=1\) je obdobně přesná formule

\[
|D_o|=8|\sin((\omega-g)/2)|\sin(g/2)
                 |\sin(\omega-\pi-g/2)|.
\tag{CG26}
\]

První prefaktor je větší než \(4/3\): z \(a>2\pi/3\) je
\((\omega-g)/2>2\pi/15>2/5\), zatímco horní konec je pod \(\pi/2\), a
\(\sin(2/5)>1/3\). Případ \(|\omega-\pi-g/2|\le10^{-3}\) je dostupný
pouze zápornému pulzu. Tentýž odhad \(|P_\omega|<6\),
\(|\kappa_\omega|<2\) by dal vzdálenost do \(8\cdot10^{-3}\) od druhé QA
čtveřice. Je tedy vyloučen. Z (CG21) je modul úhlového rozdílu v (CG26)
menší než \(\pi/2\), a \(\sin y\ge2y/\pi>(3/5)y\) dává \(|D_o|>1/1250\).

Obě formule tak poskytují jednotný původní, fyzikální odhad

\[
\boxed{|D_o|>1/2000\quad\hbox{na každém hranovém průchodu v BF pásu}.}
\tag{CG27}
\]

Použitá malá pulzní okolí nejsou odvozena z extrapolace starého
výstupního okna. U první QA skupiny, na celé spojovací úsečce
\(|\omega-g/2|\le10^{-3}\), je \(|\varepsilon-3\pi/5|\le2/1000\), tedy
\(\cos\varepsilon>-31/100-2/1000=-39/125>-1/3\) a \(W^2>2/5>1/4\). U
druhé skupiny obdobně \(|\varepsilon-2\pi/5|\le2/1000\),
\(\cos\varepsilon>307/1000>0\) a \(W^2>3/5>1/4\). Obě meze přímo
zajišťují \(|\kappa_\omega|<2\). Také \(|v|<1/20+5/1000=11/200<1/10\),
takže \(|P_\omega|\le5/\sqrt{1-v^2}<6\) na celých těchto úsecích.

\hypertarget{hlavnuxed-zuxe1pornuxe1-vuxe1zanuxe1-vux11btev-pux159i-bf-energii-pruxe1vux11b-jeden-koux159en}{%
\subsection{Hlavní záporná vázaná větev: při BF energii právě jeden
kořen}\label{hlavnuxed-zuxe1pornuxe1-vuxe1zanuxe1-vux11btev-pux159i-bf-energii-pruxe1vux11b-jeden-koux159en}}

Tato část rozlišuje počet \textbf{stacionárních} řešení od počtu všech
energetických kořenů. Druhý počet nevyplývá pouze ze slovního soupisu
stacionárních křivek; zde se odvozuje přímo.

Při pevném \(e\in I_e\) parametrizujme všechny přístupné body hlavní
záporné větve pomocí \(0\le\delta\le\delta_{\max}(e)\):

\[
p(\delta)=\arccos v(\delta),\quad P(\delta)=\pi-p(\delta),
\quad \varepsilon=d(e)+\delta,\quad d(e)=\pi+g-e,
\quad \kappa(\delta)\ge0.
\tag{CG28}
\]

Čtyři možné celkové hybnosti jsou \(\pm K_-(\delta)\),
\(\pm K_+(\delta)\), modulo \(2\pi\), kde

\[
K_-=P-\kappa,\qquad K_+=P+\kappa.
\]

V otevřeném parametrickém intervalu je

\[
K_-'=-\frac1w+\frac1W=\frac{w-W}{wW},
\qquad K_+'=-\frac1w-\frac1W<0.
\tag{CG29}
\]

BS23 má právě jednu nulu a mění se tam z minus do plus. Proto \(K_-\) má
právě jedno přísné \textbf{globální minimum}, rovné \(K_b(e)\), a
dosahuje ho pouze při \(\delta=\delta_*\). Všude jinde je \(K_->K_b\).

Na vnitřku intervalu je \(K_+>K_-\ge K_b\). Při konci \(\delta_{\max}\)
sice \(\kappa=0\), ale \(K_+=K_->K_b\), protože minimum je vnitřní. Také
na začátku jsou hodnoty od minima odděleny. Protože \(0<K_-<\pi\),
záporná kopie \(-K_-\) nemůže dát kladné \(K_b\) ani po redukci do
\((-\pi,\pi]\).

Zbývá možnost \(-K_+\) po oběhu kruhu. Platí \(0<K_+\le\pi+\kappa(d)\),
a \(\kappa(d)<\pi/2\): \(d>\pi/2\), tudíž
\(\cos\kappa(d)=(1-5\cos d)/4>0\). Případná kladná redukovaná hodnota
\(-K_+\) je proto větší než \(\pi/2>6/5>K_b\). Žádná další znaménková
kopie nepřidává kořen.

\[
\boxed{\begin{gathered}
\text{Při }(K,e)=(K_b(e),e)\text{ má hlavní záporná větev}\\
\text{právě jediný energetický kořen P: původní dvojitý BF bod.}
\end{gathered}}
\tag{CG30}
\]

Reflexe dokazuje totéž u záporného BF. Pole má původní kladnou kontaktní
váhu a křivost \(>1/4\); nebylo odstraněno ani nahrazeno algebraickým
kandidátem.

Proto je reálná spojená fáze této větve při pevném \(K_b\) všude alespoň
\(e\), s rovností jen v BF bodě. Důvodem je její spojitost na celém
kruhu, kladný lokální kvadratický rozdíl u BF a (CG30): změna znaménka
jinde by vyžadovala další energetický kořen. Tento argument sám však
nedává číselnou velikost kladné mezery mimo zvolený Morseův interval.

\hypertarget{co-tento-vuxfdsledek-duxe1vuxe1-celuxe9mu-komplementu-a-co-jeux161tux11b-potux159ebuje}{%
\subsection{Co tento výsledek dává celému komplementu a co ještě
potřebuje}\label{co-tento-vuxfdsledek-duxe1vuxe1-celuxe9mu-komplementu-a-co-jeux161tux11b-potux159ebuje}}

Pro celý původní \(F_2\), po odečtení konkrétního lokálního BF
intervalu, jsou nyní doloženy následující údaje:

\begin{enumerate}
\def\labelenumi{\arabic{enumi}.}
\tightlist
\item
  Všechny jiné fyzikální kritické konfigurace jsou od centra BF
  vyčísleně odděleny, s nejmenší zde použitou kruhovou rezervou
  \(3/100\). Platí to na celém původním kruhu, bez omezení na vybrané
  výstupní okno.
\item
  Každá reálná kontinuální hrana v BF pásu je jednoduchá:
  \(|\partial_Pt|>1/1600\), a její původní kontaktní jmenovatel je
  \(|D_o|>1/2000\).
\item
  Na hlavní záporné větvi není vedle vyjmutého BF bodu žádný další
  energetický pól. Hlavní kladná a obě nízké větve jsou nestacionární
  podle předchozích rychlostních důkazů; jejich skutečné příspěvky
  zůstávají v komplementu.
\item
  Samostatné společné příspěvky QA, QC a QB již mají odvozené lokální
  holomorfní obvody, ale zde získaná separace navíc dokazuje, že tyto
  společné konfigurace nenastávají přímo na BF centrech.
\end{enumerate}

Tyto body \textbf{neopravňují} nahradit lokální mez
\(|B_{\rm loc}|<2^{80}\) mezí pro celý \(B=B_{\rm loc}+F_{2,\rm comp}\).
Ke číselné normě celého komplementu zbývá převést doložené reálné
informace na konkrétní komplexní obvod: uvést jeho vzdálenost od všech
příslušných kořenů, minima jmenovatelů na obvodu a jednotnou mez
integrandu. Zejména v části hlavní záporné větve mimo maličký BF
interval je pozitivita energetické mezery z (CG30) kvalitativní, dokud
není zaplacen konkrétní obvod nebo přímý kvantitativní odhad této
mezery.

Jde tedy o dokončený fyzikální geometrický certifikát pro tento
zbývající krok, nikoli o nový obecný časový důkaz nebo nový čas. Závěr
původního experimentu zůstává podmíněn všemi jeho dřívějšími aparátovými
a rozptylovými předpoklady.

Konkrétní mezera posledního uvedeného atomového dílu je dále
vyčíslena v AM13--AM18. Norma celého odpovídajícího atomového
komplementu je doložena v AM3. Tyto následné výsledky nepřidávají
do CG27 žádnou dosud neodvozenou mez kontinua.

\clearpage
\section{Holomorfie celého komplementu a výslovné koncové meze}
\hypertarget{pux16fvodnuxed-pux159esnuxe1-dekompozice}{%
\subsection{Původní přesná
dekompozice}\label{pux16fvodnuxed-pux159esnuxe1-dekompozice}}

Pišme jako v BF

\[
\Phi(P,K)=\beta+\omega_b(P)-\varepsilon(K-P),\qquad
\Phi(P(x,K),K)=E_b(K)+x^2,
\] \[
Z=E-E_b(K),\qquad r=2^{-64}.
\tag{HC1}
\]

Původní místní příspěvek \(F_{2,\mathrm{loc}}\) je celý člen
\(\sigma=-1\) na intervalu \(P(-r,K)\le P\le P(r,K)\), s původní mírou
\(\mathrm dP/(2\pi)\). Jeho doplněk je přesně

\[
F_{2,\mathrm{comp}}=F_2-F_{2,\mathrm{loc}}.
\tag{HC2}
\]

Do doplňku patří celá druhá pulzní větev a celý zbytek první větve.
Pohyb konců v K se zachovává na obou stranách rovnosti. Zdroj, jeho
normalizace, celé energetické okno, kruh K, všech 26 štítků a původní
pravděpodobnostní jev se tím nemění.

\hypertarget{vuxfdslovnux11b-omezenuxe1-ux10duxe1st-u-pohyblivuxfdch-koncux16f}{%
\subsection{Výslovně omezená část u pohyblivých
konců}\label{vuxfdslovnux11b-omezenuxe1-ux10duxe1st-u-pohyblivuxfdch-koncux16f}}

Na širším Morseově oboru z BF24 jsou doloženy

\[
|K-K_0|<2^{-60},\quad |x|<2^{-61},\quad
|a(x,K)|<2^{12},\quad |P_x|<8,
\] \[
a=h_-r_bP_x,\quad |h_-|<5,\quad |R_A|<2^{34},\qquad
f_A=\frac{2ir_b}{Z-x^2}+R_A.
\tag{HC3}
\]

Vyčleňme dva přesné intervaly \(r\le|x|\le2r\). Jejich celý původní
příspěvek označme \(F_{2,\mathrm{ann}}\). Zvolme

\[
|K-K_0|<2^{-60},\qquad |Z|<r^2/4=2^{-130}.
\tag{HC4}
\]

Všechny zde používané objektové argumenty zůstávají v oboru BF18.
Skutečně \(|P-P_0|\le8|K-K_0|+4|x|<2^{-57}+2^{-61}\), zatímco
\(|\omega-\omega_b(P)|=|Z-x^2|<5r^2\). Obě meze jsou přísně menší než
odpovídající části objektového poloměru \(R_o=2^{-24}\).

Na uvedených intervalech

\[
|Z-x^2|\ge\frac34x^2.
\tag{HC5}
\]

Přesný integrál proto je

\[
F_{2,\mathrm{ann}}=
\frac{i}{\pi}\int_{r\le|x|\le2r}\frac{a(x,K)}{Z-x^2}\,\mathrm dx
+\frac1{2\pi}\int_{r\le|x|\le2r}h_-R_AP_x\,\mathrm dx.
\tag{HC6}
\]

Je holomorfní v obyčejných proměnných K,Z na celém (HC4), včetně
\(Z=0\). První člen je omezen číslem

\[
\frac{2^{12}}{\pi}\frac43\,2\int_r^{2r}\frac{\mathrm dx}{x^2}
=\frac{2^{12}}{\pi}\frac4{3r}<2^{77}.
\]

Druhý člen má modul menší než \(2r\cdot5\cdot2^{34}\cdot8<2^{-23}\).
Celkem tedy bezpečně

\[
\boxed{|F_{2,\mathrm{ann}}|<2^{78}\quad\text{na (HC4)}.}
\tag{HC7}
\]

Pro reálné K a původní horní komplexní energii jde o prostou změnu
integrační proměnné. Následné pokračování má proto přesně původní
fyzikální význam. Kolem konců \(x=\pm r\) nevzniká nezapočtené reziduum
ani nový prahový bod.

Na menším původním oboru BF2 je k dispozici Cauchyův polydisk poloměrů
\(2^{-61}\) v K a \(2^{-131}\) v Z. Pro všechna \(m,n\ge0\) tudíž

\[
\boxed{
|\partial_K^m\partial_Z^n F_{2,\mathrm{ann}}|
\le m!n!\,2^{61m+131n+78}.}
\tag{HC8}
\]

Derivace v K je zde při pevném Z; \(\partial_E=\partial_Z\). Část
\(r\le|x|\le2r\) tedy již není otevřeným číselným úkolem. Mez (HC7) se
nevztahuje na vzdálený zbytek integračního kruhu.

\hypertarget{uxfaplnost-muxedstnuxedch-typux16f-mimo-vyjmutuxfd-vuxe1zanuxfd-bod}{%
\subsection{Úplnost místních typů mimo vyjmutý vázaný
bod}\label{uxfaplnost-muxedstnuxedch-typux16f-mimo-vyjmutuxfd-vuxe1zanuxfd-bod}}

Fixujme jedno původní reálné centrum \((K_0,E_0)\) na BF. Na každém z
obou pulzních příspěvků rozdělme původní kruh na intervaly s konci mimo
všechny singularity. Nepoužívá se hladké dělení jednotky v P: takový
váhový činitel by sám neumožňoval holomorfní deformaci obvodu.

Úplný objektový vzorec je racionální v původní fyzikální odmocnině. Jeho
možné místní překážky tedy pocházejí z objektových hran, kontaktních
pólů a výměny pulzních větví. Jejich úplný fyzikální soupis již byl
dokázán v PT, BC, M, L a BS.

\begin{enumerate}
\def\labelenumi{\arabic{enumi}.}
\tightlist
\item
  \textbf{Obyčejná objektová hrana.} Jestliže je relativní rychlost
  nenulová, její průmět je jednoduchý holomorfní kořen v P. Původní
  vnější energie posune kořen do jedné poloroviny. Pevný oblouk v opačné
  polorovině dává holomorfní místní periodu. Obě stacionární kontinuální
  křivky jsou od našeho centra odděleny podle CS12.
\item
  \textbf{Jednoduchý objektový pól uvnitř mezery.} Pro nenulovou
  relativní rychlost je jeho průmět opět jednoduchý. Původní jmenovatel
  splňuje \(D_P=D_\omega(\sigma\varepsilon'-\omega_b')\ne0\). Stejná
  vnější deformace zachovává fyzikální reziduum a dává holomorfní místní
  příspěvek. Jediný stacionární pól při našem \((K_0,E_0)\) je právě
  vyjmutý BF bod. Druhá vázaná konfigurace má jiné K podle BS25--BS30.
\item
  \textbf{Pól společně s hranou.} Všechny možnosti jsou již úplně
  zahrnuty v QA. Jejich společná perioda má holomorfní pokračování, i
  když jednotlivé oddělené příspěvky takový závěr nedávají.
\item
  \textbf{Kuželový konec.} Jediný dosažitelný objektový kužel je zahrnut
  v QC, včetně obou jeho hran a původního listu.
\item
  \textbf{Výměna pulzních větví.} Přesné maticové rozštěpení KP15b
  zachovává obě původní větve. Mimo objektovou hranu je součet
  regulární; na současné hraně jej řeší QB.
\end{enumerate}

Žádný vynechaný stacionární fyzikální typ v tomto seznamu nezůstává.
Póly na jiné odmocninové vrstvě se do něj svévolně nepřidávají. Současné
kořeny dvou různých pulzních sčítanců se deformují v každém z přesně
zachovaných sčítanců zvlášť.

\hypertarget{kompaktnuxed-objektovuxe9-konce-nevytvuxe1ux159ejuxed-dalux161uxed-pux159ekuxe1ux17eku}{%
\subsection{Kompaktní objektové konce nevytvářejí další
překážku}\label{kompaktnuxed-objektovuxe9-konce-nevytvuxe1ux159ejuxed-dalux161uxed-pux159ekuxe1ux17eku}}

Je třeba ošetřit také \(v=\pm1\), kde kontaktní váha některého
objektového stavu může vymizet. Samotné zkrácení funkce při jednom P by
k holomorfii okolní periody nestačilo.

Při \(v=1\) se původní polynomy přesně rozkládají jako

\[
D_o=(x-1)N,\qquad
N=x^2+\frac35(1+q)x+q.
\tag{HC9}
\]

U viditelného kompaktního pólu \(x=1\), \(\omega=g\), je

\[
N=\frac85(1+q)\ne0,\qquad
D_\omega=-\frac{8i}{5}(1+q)\ne0.
\tag{HC10}
\]

Zde \(q\ne-1\), protože \(t=1-5\cos g\) je přísně uvnitř mezery a není
roven \(-4\). Půlperiodová symetrie dává totéž pro viditelný kompaktní
pól při \(v=-1\).

Zrušený faktor N nemá při \(v=1\) přehlédnutý kořen na otevřeném
kontinuu. Tam je fyzikální q reálné, \(|q|<1\), a při
\(x=e^{-i\delta}\), \(\delta=\omega-g\), platí přesně

\[
\frac Nx=(1+q)(\cos\delta+3/5)
-i(1-q)\sin\delta.
\]

Nula by tedy vyžadovala x=±1, ale \(N(1)=8(1+q)/5\) a \(N(-1)=2(1+q)/5\)
jsou nenulové. Symetrie řeší i \(v=-1\). Na hranách byly všechny
fyzikální nuly \(D_o\) úplně určeny v PT23: mají \(v=\pm v_a\),
\(0<v_a<1\), a proto neleží na kompaktních koncích. Všechny zbývající
koncové kořeny tak skutečně leží v přísné fyzikální mezeře.

U ostatních koncových objektových kořenů je
\(L=(\cos(\omega-g)-v)^2+16\sin^2(\omega-g)>0\). Původní reálná fázová
rovnice BC10--BC12 se tam proto hladce pokračuje a stále má
\(F_\omega>1\). Výslovně s
\(B_o=-(2/5)x[\cos\delta-v+4i\sin\delta]\ne0\), \(q=e^{-i\eta_o}\),
\(q_{\rm pole}=e^{-i\varphi}\) je \(D_o=B_o(q-q_{\rm pole})\). V jeho
nule proto
\(D_{o,\omega}=iB_oq\,\partial_\omega(\varphi-\eta_o) =iB_oq F_\omega\ne0\).
Objektová větev je holomorfní přes konec v P; její kontaktní reziduum
smí být nulové, ale nevzniká nový pól vyššího řádu.

Na kompaktních koncích je objektová rychlost nulová. Původní energetické
podmínky a úplná vyloučení M, L, BS zajišťují nenulovou pulzní rychlost
u všech zde dosažitelných rezonancí. Jde tedy opět o nestacionární
případ bodu 2 výše, případně o odstranitelný čitatel. Zkrácení nevytváří
nový stacionární příspěvek celého druhého kontaktu.

\hypertarget{sloux17eenuxed-pux16fvodnuxedch-vnux11bjux161uxedch-period}{%
\subsection{Složení původních vnějších
period}\label{sloux17eenuxed-pux16fvodnuxedch-vnux11bjux161uxedch-period}}

Po odstranění intervalu \(|x|<r\) jsou na reálném kruhu při pevném
\((K_0,E_0)\) jen konečně mnohé výše klasifikované singularity. Plyne to
také přímo z algebraického popisu původního integrandu. Identicky nulová
kontaktní větev by vytvořila celý interval stacionárních rezonancí,
který úplný soupis vylučuje.

Kolem každého zbývajícího bodu lze zvolit malý interval a správný pevný
vnější oblouk podle předchozího seznamu. Jejich konce jsou regulární. Na
zbytku kruhu je původní fyzikální integrand holomorfní v okolí kompaktní
cesty. Konečný součet těchto period je proto holomorfní na společném
komplexním okolí parametrů. Původní horní energetický limit je ve všech
členech stejný.

Případné spoje mezi souřadnicí P a podepsanou pulzní souřadnicí
\(h=k-\pi\) lze vést po regulárních úsecích. Pohyb konců s K pak přidává
pouze integrály holomorfního integrandu mezi holomorfními konci. Obdobně
jsou konce vyjmutého BF intervalu kontrolovány již výslovně v
(HC4)--(HC8).

Tím se dokazuje skutečné tvrzení pro celý původní zbytek:

\[
\boxed{
F_{2,\mathrm{comp}}\text{ má u každého původního centra vázaného prahu}
\text{ holomorfní pokračování v }(K,Z).}
\tag{HC11}
\]

Používá se přesný součet integračních úseků. Žádný zdrojový stav, pulzní
větev, kontakt ani část vybraného pravděpodobnostního jevu se
neodstraňuje.

Morseova souřadnice
\(x=\operatorname{sgn}(P-P_b) \sqrt{\Phi(P,K)-E_b(K)}\), s kladnou
orientací v reálném řezu, je na překryvech jednoznačná. Pevné konce
\(x=\pm r\) se tudíž na překryvech shodují. Obě původní BF křivky nad
uzavřeným I jsou kompaktní. Konečný výběr těchto okolí proto zaručuje
také společný kladný poloměr a konečnou horní mez zbytku po vhodném
zmenšení. \textbf{Tento poslední argument jejich číselné hodnoty
neurčuje.}

\hypertarget{dux16fsledek-pro-celuxfd-pux159enos-a-zbuxfdvajuxedcuxed-ux10duxedselnuxfd-uxfadaj}{%
\subsection{Důsledek pro celý přenos a zbývající číselný
údaj}\label{dux16fsledek-pro-celuxfd-pux159enos-a-zbuxfdvajuxedcuxed-ux10duxedselnuxfd-uxfadaj}}

S \(B=B_{\mathrm{loc}}+F_{2,\mathrm{comp}}\) a \(s^2=Z\) je původní celý
přenos přesně

\[
V(K,E_b+s^2)=
\frac{(-2A-2sB,\,4s)^T}{F_1A+s(4+F_1B)}.
\tag{HC12}
\]

Z (HC11), BF a FQ je tedy skutečně holomorfní funkcí \((K,s)\) na
nějakém okolí \(s=0\). Má rozvoj

\[
V(K,E_b+s^2)=\binom{-2/F_1(E_b)}0
+s\binom{8/[F_1(E_b)^2 A(K,0)]}{4/[F_1(E_b)A(K,0)]}
+O(s^2).
\tag{HC13}
\]

Koeficient s je nenulový. Tento kvalitativní tvar již není podmíněn
dodatečnou holomorfií celého komplementu; ta byla výše dokázána. Jeho
společný číselný obor a bodová konstanta však nadále vyžadují explicitní
společný obvod a mez celého zbytku.

Je-li taková mez na konkrétním společném oboru doložena jako
\(|B|\le M_B\), již ověřené původní FQ konstanty dovolují

\[
|K-K_0|<2^{-88},\qquad
|s|<\min\left\{2^{-76},\frac{2^{-13}}{4+3M_B}\right\},
\qquad
\left\|V-\binom{-2/F_1(E_b+s^2)}0\right\|_2<2^{21}|s|.
\tag{HC14}
\]

K číselnému oboru v (HC14) se ovšem musí přidat skutečný explicitní
obor, na němž platí dodaná mez \(M_B\). Zde přidělené \(2^{78}\) pro
koncové úseky ani starší \(2^{80}\) pro místní \(B_{\mathrm{loc}}\) se
za normu celého B nepovažují.

Samostatný nový odhad SP používá spektrální kladnost a poskytuje
oboustrannou integrovanou stavovou mez bez této bodové konstanty. Jeho
normový závěr není zaměňován s chybějící číselnou holomorfní mezí v
tomto oddílu. Žádný nový globální čas zde tvrzen není.

\clearpage
\section{Číselná mez celého hlavního záporného atomového komplementu}
\hypertarget{pux159esnuxfd-vuxfdsledek}{%
\subsection{Přesný výsledek}\label{pux159esnuxfd-vuxfdsledek}}

Fixujme kladný střed vázaného prahu \((P_0,K_0,E_0)\),
\(e_0=E_0-\beta\), \(K_0=K_b(e_0)\). Záporný střed se získá reflexí.
Označme původní hlavní zápornou objektovou fázi

\[
\omega_b(P)=\pi+\omega_c(-\cos P),\qquad
u_b(P)=e^{-i\omega_b(P)},
\tag{AM1}
\]

a její skutečnou kontaktní atomovou váhu \(r_b(P)\). Pro reálné P platí
\(0\le r_b(P)\le1\), protože jde o atom původní normalizované kontaktní
spektrální míry. Případné kompaktní koncové hodnoty lze doplnit spojitě;
hodnoty na jednotlivých bodech nemění integrál.

Nechť \(P(x,K)\), \(E_b(K)\) jsou přesně funkce BF24 a

\[
r=2^{-64},\qquad
a(K)=P(-r,K),\quad b(K)=P(r,K),\quad I(K)=[a(K),b(K)].
\tag{AM2}
\]

Při komplexním K znamená interval holomorfní pokračování téže Morseovy
cesty. Definujme \(\mathcal A_{\mathrm{comp}}\) jako celý původní
příspěvek atomu (AM1), s oběma pulzními větvemi, minus jeho
\(\sigma=-1\) příspěvek na \(I(K)\). Pak

\[
\boxed{\begin{gathered}
|K-K_0|<2^{-286},\qquad
Z=E-E_b(K),\quad |Z|<2^{-286}\\
\Longrightarrow\quad
\mathcal A_{\mathrm{comp}}\text{ je holomorfní a }
|\mathcal A_{\mathrm{comp}}|<2^{280}.
\end{gathered}}
\tag{AM3}
\]

Obor je menší než BF2. Funkce je pokračováním původního vnějšího
integrálu; žádný atom ani pulzní větev se neodstraňuje změnou
normalizace. Následuje úplný důkaz včetně pevného intervalu a obou
koncových oprav.

\hypertarget{jednotnuxe9-pomocnuxe9-meze}{%
\subsection{Jednotné pomocné meze}\label{jednotnuxe9-pomocnuxe9-meze}}

Z CG10--CG13 a BS vyplývá ve středu

\[
\frac3{10}<\delta_*<\frac{13}{40},\quad
-\frac35<v_*<-\frac25,\quad
\frac34<K_0<\frac65.
\]

Na celém BS boxu \(1/4<\delta<1/3\), \(1<z<5/4\) máme

\[
-v'=\sin\delta+4\cos\delta\,z+4\sin\delta\,z'
<\frac13+5+\frac{16}{9}=\frac{64}{9}.
\]

Protože ve středu \(\sin P_0>4/5\), jeho objektová i pulzní rychlost
splňuje

\[
\boxed{w_* >9/80.}
\tag{AM4}
\]

Na hlavní větvi je celkově \(|\omega_b'(P)|<3/4\) podle M3 a fyzikální
symetrie. Pulz má pro \(0\le k\le\pi\)

\[
W(k)=|\varepsilon'(k)|
=\frac{2\sin(k/2)}{\sqrt{3-2\cos k}},
\qquad
\boxed{k/5<W(k)\le k\quad(k>0).}
\tag{AM5}
\]

Horní mez plyne ze \(2\sin(k/2)\le k\) a jmenovatele nejméně 1. Pro
dolní použijme \(\sin(k/2)\ge k/\pi\), \(\sqrt5<3\), \(\pi<10/3\). Dále
vždy \(W<1\). Každá reálná spojená fáze
\(\omega_b(P)-\varepsilon(K_0-P)\) je tedy na kruhu Lipschitzova s
konstantou menší než 2.

Při pevném \(e_0\) pišme jako v CG28

\[
d=\pi+g-e_0,\quad \delta_{\max}=a_p-d,\quad
a_p=\arccos(-3/5),\quad
P(\delta)=\pi-\arccos v(\delta).
\]

Označme \(p_e=P(\delta_{\max})\). Protože
\(\kappa_*\ge W(\kappa_*)=w_*>9/80\), dostáváme z (AM5)

\[
\delta_{\max}-\delta_*
=a_p-\varepsilon(\kappa_*)
=\int_0^{\kappa_*}W(t)\,dt
>\frac{81}{64000}>\frac1{1000}.
\tag{AM6}
\]

Na této části BS dává \(|d\delta/dP|<1/3\), a proto \(P_0-p_e>3/1000\).

Budeme také potřebovat malou explicitní rezervu od konce BS boxu:

\[
\boxed{\delta_{\max}<1/3-1/2000.}
\tag{AM7}
\]

Je to přímé zesílení BS11. Položme \(\Delta=a_p-3\pi/5\),
\(x=997/3000\). BS dává \(\tan\Delta<343/1000\). Přitom \[
\tan x\ge x+x^3/3
>\frac{1032937}{3000000}
>\frac{343}{1000}+\frac1{1000}.
\] Na \([0,x]\) je \(\sec^2<2\), takže \(x-\Delta>1/2000\). Z
\(\delta_{\max}\le\Delta+1/1000\) plyne (AM7).

\hypertarget{kvantitativnuxed-rux16fst-energetickuxe9-kux159ivky-k-pux159i-pevnuxe9m-e0}{%
\subsection{Kvantitativní růst energetické křivky K při pevném
E0}\label{kvantitativnuxed-rux16fst-energetickuxe9-kux159ivky-k-pux159i-pevnuxe9m-e0}}

Použijme původní \(\Phi(P,K)=\beta+\omega_b(P)-\varepsilon(K-P)\) a
položme

\[
h=2^{-60},\qquad
\mathcal D(P,K)=\frac{\Phi(P,K)-E_0}{\Phi_K(P_0,K_0)}.
\tag{AM8}
\]

Ve středu je \(\mathcal D_P=0\), \(\mathcal D_K=1\) a \(\Phi_K=-w_*\).
Na komplexním polydisku poloměru h v P,K jsou z BF16 a pulzních mezí BF
druhé derivace \(\Phi\) menší než \(2^{21}\), třetí menší než
\(2^{44}\). Z (AM4) tedy druhé derivace \(\mathcal D\) mají mez
\(2^{25}\).

Na K-kružnici \(|K-K_0|=h\), pro \(|P-P_0|\le h\), máme

\[
|\mathcal D-(K-K_0)|
\le\tfrac12 2^{25}(|P-P_0|+|K-K_0|)^2
\le2^{26}h^2<h.
\tag{AM9}
\]

Rouchého věta dává jediný holomorfní graf \(K=\mathcal K(P)\). Je reálný
pro reálné P a shoduje se s původním \(K_-(\delta(P))\). Označení
k\textless0 je zachováno, protože jeho středová vzdálenost od nuly je
větší než \(9/80\). Dostáváme

\[
|\mathcal K-K_0|\le2^{26}h^2,
\qquad |\mathcal K'|<2^{27}h=2^{-33}.
\tag{AM10}
\]

Na tomto reálném grafu je \(\Phi_K<0\), \(|\Phi_K|<1\). Z
\(\Phi_{PP}(P_0,K_0)>1/4\) a meze třetích derivací plyne
\(\Phi_{PP}>1/4-2^{-15}\). Smíšené druhé derivace
\(\Phi_{PK},\Phi_{KK}\) mají pulzní mez 8. Proto

\[
\mathcal K''
=-\frac{\Phi_{PP}+2\Phi_{PK}\mathcal K'
               +\Phi_{KK}(\mathcal K')^2}{\Phi_K}
>\frac18.
\tag{AM11}
\]

V čitateli stačí odečíst \(2^{-15}+2^{-29}+2^{-63}\) od \(1/4\).
Následkem je

\[
\mathcal K(P)-K_0\ge\frac{|P-P_0|^2}{16}
\quad (|P-P_0|\le h).
\tag{AM12}
\]

CG29 dokazuje, že \(K_-(\delta)\) je přísně monotónní na obou stranách
svého jediného globálního minima. Proto se (AM12) převádí na celý
energeticky přípustný kladný interval \(p_e\le P\le\pi\):

\[
K_-(\delta(P))-K_0
\ge\frac{\min\{|P-P_0|^2,h^2\}}{16}.
\tag{AM13}
\]

Oba okraje lokálního intervalu jsou uvnitř přípustného intervalu: vlevo
to zajišťuje (AM6), vpravo \(P_0<6/5<\pi\). Není zde dosazen pouze
kvalitativní argument kompaktností.

\hypertarget{globuxe1lnuxed-reuxe1lnuxe1-energetickuxe1-mez-mimo-pux159esnuxfd-interval-i_0}{%
\subsection{\texorpdfstring{Globální reálná energetická mez mimo přesný
interval
\(I_0\)}{Globální reálná energetická mez mimo přesný interval I\_0}}\label{globuxe1lnuxed-reuxe1lnuxe1-energetickuxe1-mez-mimo-pux159esnuxfd-interval-i_0}}

Položme \(I_0=I(K_0)\). V BF Morseově souřadnici
\(x=(P-P_0)\sqrt{C_M}\), s \(0<C_M<2\) na reálném řezu, proto oba okraje
splňují

\[
|P(\pm r,K_0)-P_0|>r/2=2^{-65}.
\tag{AM14}
\]

\hypertarget{kladnuxe9-energeticky-pux159uxedpustnuxe9-p}{%
\subsubsection{Kladné energeticky přípustné
P}\label{kladnuxe9-energeticky-pux159uxedpustnuxe9-p}}

Pro \(P\in[p_e,\pi]\setminus I_0\) určuje energie \(e_0\) požadovanou
pulzní fázi \(\varepsilon_{\rm req}=d+\delta(P)\) a moment
\(\kappa_{\rm req}\). Protože \(K_-(P)=P-\kappa_{\rm req}\ge K_0\),
skutečný moment má modul \(\kappa_{\rm act}=P-K_0\ge\kappa_{\rm req}\).
Z (AM13),(AM14)

\[
\kappa_{\rm act}-\kappa_{\rm req}
=K_-(P)-K_0>2^{-134}.
\]

Skutečný energetický rozdíl tedy splňuje

\[
\begin{aligned}
g(P)&=\omega_b(P)-\varepsilon(K_0-P)-e_0\\
&=\int_{\kappa_{\rm req}}^{\kappa_{\rm act}}W(t)\,dt
>\frac{(\kappa_{\rm act}-\kappa_{\rm req})^2}{10}
>2^{-272}.
\end{aligned}
\tag{AM15}
\]

Ve zvláštním bodě \(P=p_e\) lze v (AM13) použít h, protože
\(P_0-p_e>3/1000>h\). Dostáváme silnější

\[
g(p_e)>2^{-252}=:g_0.
\tag{AM16}
\]

\hypertarget{zuxe1pornuxe9-energeticky-pux159uxedpustnuxe9-p}{%
\subsubsection{Záporné energeticky přípustné
P}\label{zuxe1pornuxe9-energeticky-pux159uxedpustnuxe9-p}}

Pišme \(P=-p\), \(p_e\le p\le\pi\). Ze CG29 máme \(p_e>K_0>3/4\).
Složený skutečný pulzní moment tedy splňuje

\[
\kappa_{\rm act}
=\operatorname{dist}_{\mathbb T}(K_0+p,0)
\ge\min\{K_0+p_e,\pi-K_0\}>3/2.
\]

Naproti tomu \(\kappa_{\rm req}\le\kappa(d)<\pi/3\): z \(\cos d<-1/5\)
plyne \(\cos\kappa(d)=(1-5\cos d)/4>1/2\). Proto
\(\kappa_{\rm act}-\kappa_{\rm req}>1/4\), a (AM5) dává \(g(-p)>1/160\).
Zejména také \(g(-p_e)>g_0\).

\hypertarget{zbyluxe9-body-s-delta-nad-energetickuxfdm-koncem}{%
\subsubsection{Zbylé body s delta nad energetickým
koncem}\label{zbyluxe9-body-s-delta-nad-energetickuxfdm-koncem}}

Pro \(|P|<p_e\) je \(\delta(P)>\delta_{\max}\), takže přímo

\[
g(P)\ge\delta(P)-\delta_{\max}.
\tag{AM17}
\]

V bodě \(p_e\) je objektová rychlost větší než \(3/32\). Skutečně BS
dává \(v_e>-13/18\), zatímco \(\delta_{\max}>\delta_*\) a monotonicita v
dávají \(v_e<v_*<-2/5\). Proto \(\sin p_e>2/3\), a \(-v'<64/9\). Na BS
boxu je navíc \(|\omega_b''|<29/48<1\), jak plyne z BF8 a jeho důkazu.

Na úsečce délky \(g_0/4\) dovnitř od \(p_e\) proto zůstává objektová
rychlost větší než \(1/16\). Tento úsek neopouští BS box: (AM7) dává
rezervu \(1/2000\), zatímco změna fáze je podle M3 menší než délka
úsečky. Pokud \(p_e-|P|\le g_0/4\), globální Lipschitzova mez 2 a (AM16)
dávají \(g(P)>g_0/2\). Pokud \(p_e-|P|\ge g_0/4\), monotonicita
objektové fáze a (AM17) dávají

\[
g(P)\ge\frac1{16}\frac{g_0}{4}=2^{-258}.
\]

Všechny případy tedy uzavírají konkrétní celokruhovou mezeru

\[
\boxed{
\omega_b(P)-\varepsilon(K_0-P)-e_0>\gamma,
\qquad\gamma=2^{-272},\qquad P\notin I_0.}
\tag{AM18}
\]

Jde o skutečnou původní fázi na celém kruhu. Převod z CG30 na číselnou
mezeru zde nevyužívá nevyčíslené minimum.

\hypertarget{celuxe1-pulznuxed-matice-pevnuxfd-reuxe1lnuxfd-komplement}{%
\subsection{Celá pulzní matice: pevný reálný
komplement}\label{celuxe1-pulznuxed-matice-pevnuxfd-reuxe1lnuxfd-komplement}}

Pro \(\lambda=e^{-ie}\) a původní pulzní symbol V definujme

\[
\boxed{
\mathcal C_P(K,e)=
(\lambda I+u_b(P)V(K-P))
\, (\lambda I-u_b(P)V(K-P))^{-1}.}
\tag{AM19}
\]

Pro reálné K a původní vnější e je \(r_b(P)[\mathcal C_P]_{11}\) přesně
celý objektový atom konvolvovaný s oběma pulzními větvemi. Tato maticová
formule nevytváří volbu pulzní banky u \(k=\pm\pi\).

Při \((K_0,e_0)\) má větev \(\sigma=-1\) podle (AM18) úhlovou spektrální
vzdálenost větší než \(\gamma\). Její rozdíl fází je menší než \(\pi\),
neboť \(\omega_b\le\pi+g+b\), \(b=\pi-a_p<\pi/3\). Pro druhou větev
platí

\[
\pi+g\le\omega_b+\varepsilon\le2\pi+g,
\]

takže její kruhová úhlová vzdálenost od \(e_0\) je větší než 1. Proto
pro \(P\notin I_0\)

\[
\|[\lambda_0I-u_b(P)V(K_0-P)]^{-1}\|<2^{273}.
\tag{AM20}
\]

Použita byla \(2\sin(\theta/2)\ge2\theta/\pi>\theta/2\) pro
\(0\le\theta\le\pi\).

Na pomocném absolutním polydisku

\[
|K-K_0|<2^{-282},\qquad |e-e_0|<2^{-282}
\tag{AM21}
\]

je \(|\lambda-\lambda_0|<2|e-e_0|\) a \(\|V(K-P)-V(K_0-P)\|<4|K-K_0|\).
Druhá mez plyne přímo z původních čtyř exponenciálních položek matice V;
její derivace má na tomto proužku Frobeniovu normu menší než 4,
nezávisle na reálném P. Neumannův součin je proto menší než
\(6\,2^{-282}2^{273}<1/2\). Získáváme

\[
\|[\lambda I-u_bV]^{-1}\|<2^{274},\qquad
\|\mathcal C_P\|<2^{277}.
\tag{AM22}
\]

Pro druhou nerovnost lze použít přesnou identitu
\(\mathcal C_P=2\lambda(\lambda I-u_bV)^{-1}-I\), \(|\lambda|<2\).

Pevný komplement

\[
J_0(K,e)=\int_{\mathbb T\setminus I_0}
r_b(P)[\mathcal C_P(K,e)]_{11}\frac{dP}{2\pi}
\tag{AM23}
\]

je tedy holomorfní na (AM21), s \(|J_0|<2^{277}\). Stačí dominované
holomorfní integrování: P zůstává reálné, \(0\le r_b(P)\le1\), a žádná
globální komplexní šířka objektové vázané fáze v P není potřebná.

Na výsledném oboru (AM3) totiž BF20 dává

\[
|e-e_0|\le8|K-K_0|+|Z|<9\,2^{-286}<2^{-282},
\tag{AM24}
\]

takže celý tento Neumannův odhad zůstává platný.

\hypertarget{pux159esnuxe9-muxedstnuxed-a-koncovuxe9-opravy}{%
\subsection{Přesné místní a koncové
opravy}\label{pux159esnuxe9-muxedstnuxed-a-koncovuxe9-opravy}}

V okolí \(I_0\) jsou obě pulzní větve analyticky oddělené. Pišme

\[
c_\sigma(P,K,e)=h_\sigma(K-P)r_b(P)
\frac{e^{-ie}+e^{-i[\omega_b(P)+\sigma\varepsilon(K-P)]}}
{e^{-ie}-e^{-i[\omega_b(P)+\sigma\varepsilon(K-P)]}}.
\tag{AM25}
\]

Na reálném P pro původní vnější e platí přesně
\(r_b[\mathcal C_P]_{11}=c_++c_-\). Položme \(a_0=a(K_0)\),
\(b_0=b(K_0)\). Požadovaný atomový komplement má přesnou reprezentaci

\[
\boxed{
\mathcal A_{\mathrm{comp}}
=J_0
 +\int_{a_0}^{b_0}c_+\frac{dP}{2\pi}
 +\int_{a_0}^{a(K)}c_-\frac{dP}{2\pi}
 -\int_{b_0}^{b(K)}c_-\frac{dP}{2\pi}.}
\tag{AM26}
\]

Nejprve se tato identita vezme pro reálné K a \(\operatorname{Im}e>0\),
kde je to prosté rozdělení původního integrálu. Poté pravá strana určuje
jeho holomorfní pokračování. Znaménka obou koncových oprav jsou dána
rozdílem \(\int_{I_0}c_- -\int_{I(K)}c_-\).

\hypertarget{duxe9lky-a-analytickuxe9-koncovuxe9-funkce}{%
\subsubsection{Délky a analytické koncové
funkce}\label{duxe9lky-a-analytickuxe9-koncovuxe9-funkce}}

BF24 dává na \(|K-K_0|<2^{-60}\), pro pevné \(x=\pm r\),

\[
|P(x,K)-P_0|\le8\,2^{-60}+4r<2^{-56}.
\]

Cauchyův odhad na polovičním K-disku tedy dává

\[
|a'(K)|,|b'(K)|<32,
\qquad
|a(K)-a_0|,|b(K)-b_0|<32|K-K_0|.
\tag{AM27}
\]

Obě krátké koncové cesty lze vzít jako přímé úsečky. Pro pevný interval
máme \(|I_0|\le8r\).

\hypertarget{pravidelnuxe1-plusovuxe1-muxedstnuxed-vux11btev}{%
\subsubsection{Pravidelná plusová místní
větev}\label{pravidelnuxe1-plusovuxe1-muxedstnuxed-vux11btev}}

Její reálná úhlová mezera je větší než 1 podle (AM19)--(AM20). Na oboru
(AM3) zůstává jmenovatel jejího skalárního resolventu větší než \(1/4\),
čitatel menší než 4. Z BF máme \(|h_+|<5\); na pevném reálném P je
\(r_b(P)\le1\). Proto

\[
\left|\int_{a_0}^{b_0}c_+\frac{dP}{2\pi}\right|
<80(8r)<1.
\tag{AM28}
\]

\hypertarget{obux11b-kritickuxe9-koncovuxe9-opravy}{%
\subsubsection{Obě kritické koncové
opravy}\label{obux11b-kritickuxe9-koncovuxe9-opravy}}

Na každém reálném základním konci je přesně
\(E_0-\Phi(a_0,K_0)=E_0-\Phi(b_0,K_0)=-r^2\). Pro P na příslušné krátké
koncové cestě dávají BF16 a pulzní meze \(|\Phi_P|<3\), \(|\Phi_K|<2\).
S (AM24),(AM27) tedy

\[
|[E-\Phi(P,K)]+r^2|
< (9+3\cdot32+2)2^{-286}
=107\,2^{-286}<2^{-279}.
\tag{AM29}
\]

Zejména \(|E-\Phi(P,K)|>r^2/2=2^{-129}\), přičemž jeho modul je menší
než \(1/2\). Pro \(|z|<1/2\) platí \(|e^{-iz}-1|\ge|z|/2\), a proto celý
skalární Cayleyův faktor v \(c_-\) má zde bezpečnou mez \(2^{133}\).
BF17 dává \(|r_b(P)|<64\) na této komplexní místní oblasti a opět
\(|h_-|<5\). Stačí tedy

\[
|c_-|<2^{144}.
\]

Součet délek obou úseček je menší než \(64\,2^{-286}\). S původní mírou
\(dP/(2\pi)\) dostáváme

\[
\left|\int_{a_0}^{a(K)}c_-\frac{dP}{2\pi}
-\int_{b_0}^{b(K)}c_-\frac{dP}{2\pi}\right|
<2^{150}2^{-286}<1.
\tag{AM30}
\]

Integrandy jsou holomorfní v P,K,E na celých obou krátkých cestách a
jejich konce závisejí holomorfně na K. Obě koncové opravy proto jsou
holomorfní, nikoli pouze omezené okrajové hodnoty.

Spojením (AM22),(AM26),(AM28),(AM30) získáváme dokonce
\(|\mathcal A_{\mathrm{comp}}|<2^{277}+2<2^{278}\). Rezerva \(2^{280}\)
v (AM3) je tedy bezpečná.

\hypertarget{co-pux159esnux11b-tato-mez-uzavuxedruxe1-a-co-ne}{%
\subsection{Co přesně tato mez uzavírá a co
ne}\label{co-pux159esnux11b-tato-mez-uzavuxedruxe1-a-co-ne}}

Rozklad (AM26) zachovává celý původní hlavní záporný objektový atom, obě
jeho pulzní větve i přesný pohyblivý interval, z něhož BF vyjímá svůj
původní \(\sigma=-1\) integrand. Zbytek plného kontaktu lze nadále psát
jako

\[
F_{2,\mathrm{comp}}
=\mathcal A_{\mathrm{comp}}
 +\bigl[\text{celé ostatní objektové spektrum}
       -\text{jeho původní místní }\sigma=-1\text{ díl}\bigr].
\tag{AM31}
\]

Druhý člen obsahuje objektové kontinuum a ostatní vázané větve. Jejich
možné jednoduché reálné průchody vyžadují vlastní fyzikální deformace a
číselné meze; (AM3) se za jejich normu nevydává. Zvláště se zde netvrdí
celý \(M_{\mathrm{comp}}\), nový globální Fourierův chvost ani nový čas.

Už doložený přínos je konkrétní: původně pouze kvalitativní mezera CG30
mimo maličký BF interval je nyní nahrazena celokruhovou mezerou
\(2^{-272}\), a celý odpovídající atomový komplement má přesné komplexní
okolí a číselnou normu (AM3).

Zdroj, pevný počátek času, vybraný jev, normalizace, všech 26 štítků a
korelace s pamětí a referencí se tím nemění.

\clearpage
\section{Časový závěr, kontrola a přesně zbývající údaje}

\subsection{Dokončený oboustranný normový krok}

Předchozí BF důkaz dával bodovou mez skutečného celého přenosu
jen na reálné straně \(Z>0\), bez znalosti komplementu.
Nové SP21 a SP28 řeší obě strany současně přímo ve fyzikálně
vážené normě. Není třeba předem omezit bodovou hustotu
\(\rho_2\) nebo celý \(B_{\rm comp}\).
Pomocná aproximace CT6 obsahuje všechny původní kanály.

Přesná cena je při uvedených parametrech
\[
2\sqrt{13/35}\,2^{-34}+2^{-29}<2\cdot10^{-9}.
\]
To současně platí pro všech 26 štítků a libovolnou původní
korelaci s pamětí a referencí. Není zde součet 26 samostatných
chyb: celý operátor působí na těchto faktorech totožně.

\subsection{Co zbývá pro bodový číselný komplement}

HC11 dokazuje kvalitativní holomorfii celého fyzikálního zbytku.
HC7--HC8 vyčíslují přesné koncové úseky a AM3 celý hlavní
záporný atom po původním místním odpočtu. Pro číselnou normu
\(F_{2,\rm comp}\) zbývá celý druhý člen AM31:
objektové kontinuum a ostatní vázané atomy, se stejným místním
odpočtem. Jejich reálné průchody jsou jednoduché a v CG mají
vyčíslené rychlosti, ale ty samy ještě neurčují komplexní
normu celé periody.

Přímá obvodová cesta musí dodat dvě společné číselné nerovnosti.
První vyloučí nuly příslušných objektových radikandů během
celé fyzikální deformace a omezí je na konečném obvodu.
Druhá totéž provede pro původní kontaktní jmenovatel \(D_o\).
Musí pokrývat celý zbývající integrační kruh, nejen malé
oblouky u již určených reálných kořenů. Původní listy, obě
pulzní složky a pohyblivé konce musejí zůstat zachovány.

Žádný navržený malý posun obvodu není v tomto dodatku
vydáván za dokončenou hodnotu \(M_{\rm comp}\).
Na druhé straně SP již tuto bodovou hodnotu uvnitř placených
prahových pásků nepotřebuje. Jde o dvě rozdílné analytické cesty.

\subsection{Jeden přesný dodatečný údaj pro nový globální čas}

Časová cesta přes pomocný stav potřebuje číselnou horní mez
\(B\) pro úplnou normu CT10, včetně přechodů CT2, původních
konců I a zbývajících skutečných kontinuálních prahů.
Samotná konečnost nebo omezenost profilu CT9 jí není.
Dodání B by již bez dalšího asymptotického kroku určilo
celočíselný čas CT13. Do výsledného experimentu se pak
přičte původní příchozí chyba, kalibrace přijímače,
konečný zdroj a konečné čtení podle [S2].

Nejde o tvrzení, že tyto údaje musí být novými fyzikálními
předpoklady. Jsou to konkrétní ještě nevyčíslené analytické
vlastnosti již zadaného původního operátoru.

\subsection{Dosavadní samostatný konečný čas}

Předchozí obecný pracovní návrh [S4] zůstává beze změny.
Za svých výslovných převzatých operátorových a aparátových
předpokladů uvádí
\[
N_-=10^{16},\qquad N_+=2^{81\,000\,000\,000\,000},
\]
\[
\varepsilon_-(N_-)+\varepsilon_+(N_+)
<2.310002\cdot10^{-7}<10^{-6}.
\]
S původními \(\varepsilon_S=\varepsilon_D=10^{-6}\)
je jeho dřívější konzervativní závěr
\[
p_{\rm finite}>3.3362176\cdot10^{-9}.
\]
Tato čísla jsou převzata z [S4]; nový dodatek je nezostřuje.
[S4] je interně přezkoumaný samostatný návrh této relace,
nikoli veřejně přijatý důkaz v PR~\#1817. Počet makrokroků
nemá v modelu doložený převod na fyzikální sekundy.
Matematický konečný protokol není uskutečněný experiment.

\subsection{Zachování původního pravděpodobnostního jevu}

Původní asymptotická pravděpodobnost z [S3] a původní projektor
Q se nepřepisují. Z normové chyby úplného stavu se pro tentýž
jev používá amplitudová nerovnost
\[
\left|\sqrt{p_{\rm finite}}-\sqrt{P_0}\right|
\le\varepsilon_S+\varepsilon_D+\varepsilon_-+\varepsilon_+.
\]
Samotná nová aproximační cena CT7 nenahrazuje pravou stranu
této úplné nerovnosti. Z CT7 se proto nevyhlašuje nová
pravděpodobnost hotového experimentu.

\subsection{Vnitřní matematický přezkum}

Každý nový modul prošel samostatným analytickým čtením v této
spolupracující relaci. Kontrola zahrnovala původní faktory
\(2\pi\), skutečné spektrální váhy, oba kraje energie,
všechny znaménkové kopie geometrických křivek a oba konce
pohyblivého intervalu. Dílčí poznámky z přezkumu byly
zapracovány do úplného textu: zejména skryté zrušené kořeny
u kompaktních konců, číselné rezervy v malých QA úsecích,
definice \(\Theta\) přes konce I a přesné teleskopování Cookovy řady.

Není to externí nezávislé přijetí výsledku, formální veřejný
replay ani nový výpočet CAS. V tomto pokračování nebyl proveden
vědecký numerický běh. Sazba dokumentu a kontrola jeho zobrazení
nejsou vědeckým důkazem.

\subsection{Veřejná návaznost}

Před začátkem tohoto pokračování byla 29.~10.~2026
v~10:28:43--10:31:14 UTC znovu přečtena veřejná práce.
PR~\#1817 zůstával otevřeným draftem na
\nolinkurl{66ef8220df3ad5764ca941b84223a9f46cf5411f},
PR~\#1815 na
\nolinkurl{32306016908c0097650277aa416fc0152b2f81a4}.
Rezervace \#1816 byla otevřená a nezměněná, bez komentářů.
Od předchozí kontroly v~10:18:13 UTC nebyly aktualizovány
žádné issues ani PR. Všech 322 vzdálených hlav souhlasilo
s předchozím úplným manifestem.

Závěrečná kontrola 29.~10.~2026 v~10:59:55--10:59:56 UTC
znovu ověřila oba PR, rezervaci, hlavní větev a STATUS.
Od 10:31:14 UTC nebyla aktualizována žádná issue ani PR.
Všech 322 vzdálených hlav zůstalo přesně shodných:
žádná přidaná, odstraněná ani změněná větev.

Main zůstával na
\nolinkurl{a193e389fb1ddb0bdc89165b015c95a17eb1e425}.
STATUS označoval ACTIVE Public Canon v106,
s obsahovým commitem
\nolinkurl{9242db4f040d67d9417eacd06612230b74248cff}.
Tato poznámka nepřebírá rezervaci \#1816, neprovádí
veřejnou změnu a nemění Canon.

\clearpage
\section{Zdroje a předchozí pracovní důkazy}

\begin{description}[style=nextline,leftmargin=1.2cm]
\item[S1]
\href{https://github.com/mathorn1973/twist-j/blob/66ef8220df3ad5764ca941b84223a9f46cf5411f/notes/C-BOUND-PAIR-RESOLVENT-AUDIT-N/PROOF.md}
{Původní algebraický audit úplných kontaktních resolventů, PR \#1817.}
Pin \nolinkurl{66ef8220df3ad5764ca941b84223a9f46cf5411f}.
Zejména §§9--10 a §§15.2--15.4.

\item[S2]
\href{https://github.com/mathorn1973/twist-j/blob/66ef8220df3ad5764ca941b84223a9f46cf5411f/notes/C-BOUND-PAIR-RESOLVENT-AUDIT-N/FINITE_PROTOCOL.md}
{Původní FINITE\_PROTOCOL, tentýž pin PR \#1817.}
Původní zdroj, celý jev, konečný aparát a úplný časový účet.

\item[S3]
\href{https://github.com/mathorn1973/twist-j/blob/32306016908c0097650277aa416fc0152b2f81a4/notes/C-BOUND-PAIR-ABSOLUTE-CAPTURE-N/PROOF.md}
{Úplný absolutní důkaz pravděpodobnosti, PR \#1815.}
Pin \nolinkurl{32306016908c0097650277aa416fc0152b2f81a4}.

\item[S4]
\href{twist_j_konecny_cas.tex}
{Konečný čas rozptylu původního paketu TWIST-J.}
Samostatný pracovní analytický návrh této relace, 29.~10.~2026,
33 stran. NON-CANONICAL/L1.

\item[S5]
\href{twist_j_vazane_vetve.tex}
{Vázané větve a úplný fyzikální soupis.}
Samostatný pracovní analytický dodatek, 29.~10.~2026,
28 stran. NON-CANONICAL/L1.

\item[S6]
\href{twist_j_ciselne_prahy.tex}
{Číselná okolí skutečných rozptylových prahů TWIST-J.}
Samostatný pracovní analytický dodatek, 29.~10.~2026,
27 stran. NON-CANONICAL/L1. Obsahuje QA, QB, QC, FQ a BF.

\item[S7]
\href{https://github.com/mathorn1973/twist-j/issues/1816}{Rezervace \#1816.}
Veřejné rozhraní běžící souběžné práce; tento dodatek je oddělený.
\end{description}

\end{document}
